\subsection{向量空间的子空间的维数与余维数}

命题4. 向量空间E的每个子空间F都是E的一个直因子，并且

\[
\dim \mathbf {F} + \dim (\mathbf {E} / \mathbf {F}) = \dim \mathbf {E}.\tag{3}
\]

由于商向量空间 E/F 是自由模，我们知道（§ 1，第 11 号，命题 21）F 是 E 的一个直因子；于是关系式 (3) 是公式 (1)（第 2 号）的一个特例。

推论 1. 若 E、F、G 是域 K 上的向量空间，则每一个线性映射的正合序列 \(0 \rightarrow E \rightarrow F \rightarrow G \rightarrow 0\) 分裂。

这是表达命题4（§1，第9号）的另一种方式。

推论2. 设 \((\mathbf{E}_i)_{0\leqslant i\leqslant n}\) 是域 K 上向量空间的一个有限族。若存在线性映射的正合序列

\[
0 \longrightarrow \mathrm{E} _ {0} \stackrel {u _ {0}} {\longrightarrow} \mathrm{E} _ {1} \stackrel {u _ {1}} {\longrightarrow} \mathrm{E} _ {2} \longrightarrow \dots \longrightarrow \mathrm{E} _ {n - 1} \stackrel {u _ {n - 1}} {\longrightarrow} \mathrm{E} _ {n} \stackrel {u _ {n}} {\longrightarrow} 0\tag{4}
\]

则关系

\[
\sum_ {2 k + 1 \leqslant n} \dim \mathrm{E} _ {2 k + 1} = \sum_ {2 k \leqslant n} \dim \mathrm{E} _ {2 k}\tag{5}
\]

成立，或者，若所有空间都是有限维的，

\[
\sum_ {i = 1} ^ {n} (- 1) ^ {i} \dim \mathrm{E} _ {i} = 0.\tag{6}
\]

设 \(\mathbf{I}_k = \operatorname{Im} u_k = \operatorname{Ker} u_{k+1}\) 对于 \(0 \leqslant k \leqslant n-1\) ; \(\mathbf{I}_{k+1}\) 因此同构于 \(\mathbf{E}_{k+1}/\mathbf{I}_k\) ，从而（公式（3）） \(\dim \mathbf{I}_k + \dim \mathbf{I}_{k+1} = \dim \mathbf{E}_{k+1}\) 对于 \(0 \leqslant k \leqslant n-2\) 并且此外 \(\dim \mathbf{I}_0 = \dim \mathbf{E}_0\) 和 \(\mathbf{I}_{n-1} = \mathbf{E}_n\) ，因此 \(\dim \mathbf{I}_{n-1} = \dim \mathbf{E}_n\) 。将 \(\dim \mathbf{E}_i\) 用它作为下列和式

\[
\sum_ {k = 0} ^ {n - 1}
\]

dim \(I_{k}\) 的函数所成的表达式代入（5）的两边，我们在每边都得到 \(\sum_{k=0}^{n-1}\dim I_{k}\) ，由此得出推论。

推论3. 若M和N是向量空间E的两个子空间，则

\[
\dim (\mathbf {M} + \mathbf {N}) + \dim (\mathbf {M} \cap \mathbf {N}) = \dim \mathbf {M} + \dim \mathbf {N}.\tag{7}
\]

只需将推论2应用于正合序列即可。

\[
0 \rightarrow \mathbf {M} \cap \mathbf {N} \rightarrow \mathbf {M} \oplus \mathbf {N} \rightarrow \mathbf {M} + \mathbf {N} \rightarrow 0
\]

（§1，第8号，命题10）考虑到以下事实：

\[
\dim (\mathbf {M} \oplus \mathbf {N}) = \dim \mathbf {M} + \dim \mathbf {N}
\]

(第2号，命题2)。

推论4. 设 \(\mathbf{F}\) 是向量空间 \(\mathbf{E}\) 的一个子空间；则 \(\dim \mathbf{F} \leqslant \dim \mathbf{E}\) ；若 \(\mathbf{E}\) 是有限维的，则关系 \(\dim \mathbf{F} = \dim \mathbf{E}\) 等价于 \(\mathbf{F} = \mathbf{E}\) .

第一个断言由(3)显然成立；进一步，若dim E有限，则关系dim F = dim E蕴含dim(E/F) = 0（由(3)），而维数为0的向量空间约化为0。

推论5. 若向量空间 \(\mathbf{E}\) 是向量子空间的一个族 \((\mathbf{F}_{\iota})\) 之和，则

\[
\dim \mathrm{E} \leqslant \sum_ {\iota} \dim \mathrm{F} _ {\iota}.\tag{8}
\]

若进一步 dim E 有限，则(8)式两边相等当且仅当 E 是族 \((\mathbf{F}_{\iota})\) 的直和.

不等式(8)由(3)以及 E 同构于 \(\bigoplus_{\iota} F_{\iota}\) 的一个商这一事实得出（§ 1，第 7 号，公式 (28)）。第二个断言是 § 1，第 10 号，命题 16 的推论 5 的一个特例，因为 (8) 两边的相等意味着除有限多个指标外 \(\dim F_{\iota} = 0\) 。

定义2. 给定向量空间E，E的子空间F的余维数（相对于E），记为 \(\operatorname{codim}_{E} F\) ，或简记为 codim F，即 E/F 的维数（等于 F 在 E 中任一补空间的维数）。

于是关系式 (3) 可写成

\[
\dim \mathrm{F} + \operatorname{codim} \mathrm{F} = \dim \mathrm{E}.\tag{9}
\]

命题5. 设 \(\mathbf{F}\) , \(\mathbf{F}'\) 是向量空间 \(\mathbf{E}\) 的两个子空间，且满足 \(\mathbf{F} \subset \mathbf{F}'\) 。则 \(\operatorname{codim}_{\mathbf{E}} \mathbf{F}' \leqslant \operatorname{codim}_{\mathbf{E}} \mathbf{F} \leqslant \dim \mathbf{E}\) 。若 \(\operatorname{codim}_{\mathbf{E}} \mathbf{F}\) 是有限的，关系式 \(\operatorname{codim}_{\mathbf{E}} \mathbf{F}' = \operatorname{codim}_{\mathbf{E}} \mathbf{F}\) 蕴含 \(\mathbf{F} = \mathbf{F}'\) .

不等式 \(\operatorname{codim}_{\mathbf{E}} \mathbf{F} \leqslant \dim \mathbf{E}\) 由(9)式显然可得；且若 \(\dim \mathbf{E}\) 有限，则关系式 \(\operatorname{codim}_{E} F = \dim E\) 蕴含 \(\dim F = 0\) ，从而 \(F = \{0\}\) 。命题的其余部分由此得出，因为

\[
\operatorname{codim} _ {\mathbf {E}} \mathbf {F} ^ {\prime} = \operatorname{codim} _ {\mathbf {E} / \mathbf {F}} (\mathbf {F} ^ {\prime} / \mathbf {F}),
\]

因为 \(\mathbf{E} / \mathbf{F}'\) 典范同构于 \((\mathbf{E} / \mathbf{F}) / (\mathbf{F}' / \mathbf{F})\) （I，第4节，第7号，定理4）。

命题6. 若M和N是向量空间E的两个子空间，则

\[
\operatorname{codim} (\mathbf {M} + \mathbf {N}) + \operatorname{codim} (\mathbf {M} \cap \mathbf {N}) = \operatorname{codim} \mathbf {M} + \operatorname{codim} \mathbf {N}.\tag{10}
\]

只需将命题4的推论2应用于正合序列即可。

\[
0 \rightarrow \mathrm{E} / (\mathrm{M} \cap \mathrm{N}) \rightarrow (\mathrm{E} / \mathrm{M}) \oplus (\mathrm{E} / \mathrm{N}) \rightarrow \mathrm{E} / (\mathrm{M} + \mathrm{N}) \rightarrow 0
\]

(§ 1，第 7 号，命题 10)并使用第 2 号，命题 2。

注意，如果 \(\mathbf{E}\) 是有限维的，(10) 是 (7) 和 (9) 的推论。

命题7. 若 \((\mathbf{F}_i)\) 是向量空间 \(\mathbf{E}\) 的子空间的一个有限族，则 \(\operatorname{codim}\left(\bigcap_i \mathbf{F}_i\right) \leqslant \sum_i \operatorname{codim} \mathbf{F}_i\) .

若 \(\mathbf{F} = \bigcap_{i}\mathbf{F}_{i}\) ，则 \(\mathbf{E} / \mathbf{F}\) 同构于诸 \(\mathbf{E} / \mathbf{F}_i\) 的直和的一个子空间（§1，第7号，公式（27））。

向量空间 E 的维数为 1（相应地，维数为 2）的向量子空间通常称为过 0 的直线（相应地，过 0 的平面）（若不会引起混淆，则简称为直线（相应地，平面）（参见第 9 节第 3 号）），这是类比经典几何的语言；E 的子空间若其余维数为 1，则称为过 0 的超平面（或简称为超平面）。超平面也可以定义为集合 S 的极大元，其中 S 是按包含关系排序的、E 中不同于 E 的向量子空间全体。包含子空间 H 的 E 的子空间与 E/H 的子空间之间存在一一对应（I，第 4 节，第 7 号，定理 4）；若 E 的维数 ≥1，则 S 非空，且说 H 是 S 中的极大元意味着 E/H 不包含不同于 \{0\} 和 E/H 的子空间，这蕴含 E/H 由它的任意非零元生成，换言之，它的维数为 1。

在维数为 \(n \geqslant 1\) 的有限维向量空间中，由公式 (3)，超平面就是维数为 n - 1 的那些子空间。

命题 8. 在域 K 上的向量空间 E 中，每个向量子空间 F 都是包含它的那些超平面的交集。

只需证明对任意 \(x \notin F\) ，存在一个包含 F 但不包含 x 的超平面 H。由假设 \(F \cap Kx = \{0\}\) ，故 F 与 Kx 的和 M 是直和。设 N 是 M 在 E 中的补子空间；则 E 是 \(H = F + N\) 与 Kx 的直和，因此 H 就是具有所需性质的超平面。

注。本节中对于向量空间的子空间所证明的大部分性质，对于维数（§ 2，第 7 号，注 2）已定义的自由模的子模并不成立。*例如，交换环的理想不一定有基，因为存在某些整环 A，其中某些理想不是主理想（VII，§ 1，第 1 号），且此类环中任意两个元素都是线性相关的（§ 1，第 11 号，注 1）。*自由 A-模 E 的子模可以是自由的，不同于 E，且与 E 具有相同的维数，正如整环 A 中的主理想所示；同样的例子还表明，自由 A-模的自由子模不一定有补子模。

\subsection{线性映射的秩}

定义3. 设E、F是域K上的两个向量空间。对于E到F的每一个线性映射u，F的子空间u(E)的维数称为u的秩，记为rg(u)。

若 \(\mathbf{N} = \operatorname{Ker}(u)\) ，则 \(\mathbf{E} / \mathbf{N}\) 同构于 \(u(\mathbf{E})\) ，由此得关系式

\[
\operatorname{rg} (u) = \operatorname{codim} _ {\mathrm{E}} (\operatorname{Ker} (u))\tag{11}
\]

，从而

\[
\operatorname{rg} (u) + \dim (\operatorname{Ker} (u) = \dim E.\tag{12}
\]

此外，由公式(3)

\[
\operatorname{rg} (u) + \dim (\operatorname{Coker} (u)) = \dim F.\tag{13}
\]

命题9. 设E、F是域K上的两个向量空间，且u:E→F是一个线性映射。

\begin{enumerate}
\def\labelenumi{(\roman{enumi})}
\item
  \(\operatorname{rg}(u) \leqslant \inf(\dim E, \dim F)\) .
\item
  假设 \(\mathbf{E}\) 是有限维的；为了使 \(\operatorname{rg}(u) = \dim \mathbf{E}\) ，必要且充分的条件是 \(u\) 是单射的。
\item
  假设 \(\mathbf{F}\) 是有限维的；为了使 \(\operatorname{rg}(u) = \dim \mathbf{F}\) ，必要且充分的条件是 \(u\) 是满射。
\end{enumerate}

这直接由关系式(12)和(13)得出。

推论。设 E 是维数为 n 的有限维向量空间，u 是 E 的一个自同态。以下性质是等价的：

\begin{enumerate}
\def\labelenumi{(\alph{enumi})}
\item
  \(u\) 是双射；
\item
  \(u\) 是单射；
\item
  \(u\) 是满射；
\item
  \(u\) 是右可逆的；
\item
  \(u\) 是左可逆的；
\item
  \(u\) 的秩为 \(n\) .
\end{enumerate}

若 E 是无限维向量空间，则存在 E 的单射（相应地，满射）自同态，它们不是双射（习题 9）。

设 K， \(K'\) 为两个域， \(\sigma:K\to K'\) 为 K 到 \(K'\) 上的一个同构，E 为一个向量 K-空间， \(E'\) 为一个向量 \(K'\) -空间，而 \(u:K\to K'\) 为关于 \(\sigma\) 的半线性映射（§1，第13号）； \(u(\mathbf{E})\) 作为 \(E'\) 的子空间，其维数也称为 u 的秩。它也是将 u 视为 E 到 \(\sigma_{*}(\mathbf{E}')\) 的线性映射时 u 的秩，因为 \(u(\mathbf{E})\) 的每一组基也是 \(\sigma_{*}(u(\mathbf{E}))\) 的基.

\subsection{向量空间的对偶}

定理4. 对偶空间 \(\mathbf{E}^*\) （对应向量空间 \(\mathbf{E}\) ）的维数至少等于 \(\mathbf{E}\) 的维数。要使 \(\mathbf{E}^*\) 为有限维，必要且充分条件是 \(\mathbf{E}\) 为有限维，且此时 \(\dim \mathbf{E}^* = \dim \mathbf{E}\) .

如果 K 是 E 的标量域，则 E 同构于某个空间 \(\mathbf{K}_{s}^{\mathrm{(I)}}\) ，从而 \(E^{*}\) 同构于 \(K_{d}^{I}\) （§2，第6号，命题10）。由于 \(\mathbf{K}_{s}^{\mathrm{(I)}}\) 是 \(K_{d}^{I}\) 的一个子空间，故 \(\dim E = \text{Card}(I) \leqslant \dim E^{*}\) （第3号，命题4的推论4）；进一步，若 I 是有限的，则 \(\mathbf{K}_{d}^{I} = \mathbf{K}_{d}^{\mathrm{(I)}}\) （参见习题3(d)）。

推论。对于向量空间 \(\mathbf{E}\) ，关系 \(\mathbf{E} = \{0\}\) 和 \(\mathbf{E}^{*} = \{0\}\) 是等价的。

定理5. 给定向量空间（在同一域 K 上）的两个正合序列及线性映射

\[
0 \rightarrow \mathrm{E} ^ {\prime} \rightarrow \mathrm{E} \rightarrow \mathrm{E} ^ {\prime \prime} \rightarrow 0
\]

\[
0 \to \mathbf {F} ^ {\prime} \to \mathbf {F} \to \mathbf {F} ^ {\prime \prime} \to 0
\]

以及 K 上的两个向量空间 G、H，则相应的序列

\[
0 \rightarrow \operatorname{Hom} (\mathrm{E} ^ {\prime \prime}, \mathrm{G}) \rightarrow \operatorname{Hom} (\mathrm{E}, \mathrm{G}) \rightarrow \operatorname{Hom} (\mathrm{E} ^ {\prime}, \mathrm{G}) \rightarrow 0
\]

\[
0 \rightarrow \operatorname{Hom} (\mathrm{H}, \mathrm{F} ^ {\prime}) \rightarrow \operatorname{Hom} (\mathrm{H}, \mathrm{F}) \rightarrow \operatorname{Hom} (\mathrm{H}, \mathrm{F} ^ {\prime \prime}) \rightarrow 0
\]

是正合的且分裂的。

这由每个向量子空间都是直因子（第3号，命题4）这一事实以及§2第1号命题1和命题2得出。

推论。对于每个正合序列

\[
0 \longrightarrow \mathrm{E} ^ {\prime} \stackrel {u} {\longrightarrow} \mathrm{E} \stackrel {v} {\longrightarrow} \mathrm{E} ^ {\prime \prime} \longrightarrow 0
\]

（即同一域 K 上向量空间与线性映射的），则序列

\[
0 \longrightarrow \mathrm{E} ^ {\prime \prime *} \stackrel {t _ {v}} {\longrightarrow} \mathrm{E} ^ {*} \stackrel {t _ {u}} {\longrightarrow} \mathrm{E} ^ {\prime *} \longrightarrow 0
\]

是正合的且分裂的。

特别地，由此可知，对于 E 的每个向量子空间 M，典范同态 \(E^{*}/M^{\prime} \rightarrow M^{*}\) （其中 \(M^{\prime}\) 是 \(E^{*}\) 中与 M 正交的子空间（ \(\S 2\) ，第 4 号））是双射的。

定理 6. 对于域 K 上的每个向量空间 E，典范映射 \(c_{\mathbf{E}}: \mathbf{E} \to \mathbf{E}^{**}\) （§2，第 7 号）是单射的；要使它为双射，当且仅当 E 是有限维的。

第一个断言以及"若 E 是有限维的则 \(c_{E}\) 是双射"这一事实，都是§2第7号命题14的特殊情形。设 E 是无限维的，于是我们可以假定 \(\mathbf{E} = \mathbf{K}_{s}^{(L)}\) ，其中 L 是无限集，从而 \(E^{*} = K_{d}^{L}\) 。设 \((e_{\lambda})_{\lambda \in L}\) 是 E 的典范基，并令 \((e_{\lambda}^{*})_{\lambda \in L}\) 为 \(E^{*}\) 中相应的坐标线性型族（§2，第6号）； \(E^{*}\) 中由诸 \(e_{\lambda}^{*}\) 生成的向量子空间正是直和 \(F' = K_{d}^{(L)}\) ，而 L 为无限的假设蕴含 \(F' \neq E^{*}\) 。于是存在一个超平面 \(H'\) 于 \(E^{*}\) 中，包含 \(F'\) （第3号，命题8），且由于 \(E^{*}/H'\) 非零，其对偶亦非零（定理4的推论），该对偶等同于 \(H''\) ，即 \(H'\) 在 \(E^{**}\) 中的正交 （§2，第6号，命题9的推论）。但 \(H'' \cap c_{E}(E)\) 包含于按 \(c_{E}\) 取的像，即 \(F'\) 在 E 中的正交的像，而它按定义为 0；因此 \(c_{E}(E) = E^{**}\) 是不可能的。

通常将把 E 等同于 \(E^{**}\) 中 \(c_{E}\) 的像所成的子空间.

设 \(\mathbf{E}, \mathbf{F}\) 是域 \(\mathbf{K}\) 上的两个向量空间， \(u: \mathbf{E} \to \mathbf{F}\) 是一个线性映射。我们将定义典范同构：

\begin{enumerate}
\def\labelenumi{(\arabic{enumi})}
\item
  从 \(\operatorname{Im}(u) = u(\mathbf{E})\) 的对偶到 \(\operatorname{Im}(^t u) = {}^t u(\mathbf{F}^*)\) 上的典范同构。
\item
  从 \(\operatorname{Ker}(u) = \overline{u}^1(0)\) 的对偶到 \(\operatorname{Coker}(^t u) = \mathrm{E}^* / {}^t u(\mathrm{F}^*)\) 上的典范同构。
\item
  从 \(\operatorname{Coker}(u) = \mathbf{F} / u(\mathbf{E})\) 的对偶到 \(\operatorname{Ker}(^t u) = {}^t u^{-1}(0)\) 上的典范同构。
\end{enumerate}

我们写 \(\mathbf{I} = \operatorname{Im}(u),\mathbf{N} = \operatorname{Ker}(u),\mathbf{C} = \operatorname{Coker}(u)\) ；由正合序列

\[
0 \longrightarrow \mathrm{N} \longrightarrow \mathrm{E} \stackrel {p} {\longrightarrow} \mathrm{I} \longrightarrow 0, \quad 0 \longrightarrow \mathrm{I} \stackrel {j} {\longrightarrow} \mathrm{F} \longrightarrow \mathrm{C} \longrightarrow 0 \tag {14}
\]

我们通过转置（定理5的推论）推导出正合序列

\[
\begin{array}{l} 0 \longrightarrow \mathrm{I} ^ {*} \stackrel {{t _ {p}}} {{\longrightarrow}} \mathrm{E} ^ {*} \longrightarrow \mathrm{N} ^ {*} \longrightarrow 0, \\ 0 \longrightarrow \mathrm{C} ^ {*} \longrightarrow \mathrm{F} ^ {*} \stackrel {{t _ {j}}} {{\longrightarrow}} \mathrm{I} ^ {*} \longrightarrow 0. \end{array}\tag{15}
\]

此外，由于 \(u = j \circ p\) , \(^{t}u = ^{t}p \circ ^{t}j\) ；因此正合序列（15）定义了 \(C^*\) 到 \(\operatorname{Ker}(^{t}u)\) 上、 \(I^*\) 到 \(\operatorname{Im}(^{t}u)\) 上以及 \(N^*\) 到 \(\operatorname{Coker}(^{t}u)\) 上的典范同构，因为 \(^{t}p\) 是单射且 \(^{t}j\) 是满射。确切地说，设 \(y \in \operatorname{Im}(u)\) , \(z \in \operatorname{Ker}(u)\) , \(t \in \operatorname{Coker}(u)\) , \(y' \in \operatorname{Im}(^{t}u)\) , \(z' \in \operatorname{Coker}(^{t}u)\) , \(t' \in \operatorname{Ker}(^{t}u)\) ；当 \(y'\) , \(z'\) , \(t'\) 分别典范地等同于 \(\operatorname{Im}(u)\) , \(\operatorname{Ker}(u)\) 和 \(\operatorname{Coker}(u)\) 上的线性形式时，则

\[
\langle y, y ^ {\prime} \rangle = \langle x, y ^ {\prime} \rangle \quad \text { for   all } x \in \mathrm{E} \text { such   that } u (x) = y;\tag{16}
\]

\begin{enumerate}
\def\labelenumi{(\arabic{enumi})}
\setcounter{enumi}{16}
\item
  \(\langle z, z' \rangle = \langle z, x^* \rangle\) ，其中 \(x^* \in \mathbf{E}^*\) 取遍所有模 \(^t u(\mathbf{F})\) 的类等于 \(z'\) 的元素;
\item
  \(\langle t, t' \rangle = \langle s, t' \rangle\) ，其中 \(s \in \mathbf{F}\) 取遍所有模 \(u(\mathbf{E})\) 的类等于 \(t\) 的元素.
\end{enumerate}

特别地，我们从这些结果推导出：

命题10. 设E、F是同一域K上的两个向量空间，且 \(u: \mathbf{E} \to \mathbf{F}\) 是一个线性映射。

\begin{enumerate}
\def\labelenumi{(\roman{enumi})}
\item
  对于 \(u\) 为单射（相应地，满射），其必要且充分条件是 \(^{t}u\) 是满射（相应地，单射）。
\item
  \(\operatorname{rg}(u) \leqslant \operatorname{rg}(^t u)\) ，并且关系 \(\operatorname{rg}(u) = \operatorname{rg}(^t u)\) 在 \(\operatorname{rg}(u)\) 有限时成立。
\end{enumerate}

第二个断言由上述内容及定理4得出。

定理7. 设E是域K上的向量空间，F是E的子空间，F'是F在E*中的正交补。

\begin{enumerate}
\def\labelenumi{(\roman{enumi})}
\item
  \(\dim \mathbf{F}' \geqslant \operatorname{codim}_{\mathbb{E}} \mathbf{F}\) ；对于 \(\dim \mathbf{F}'\) 为有限，必要且充分条件是 \(\operatorname{codim}_{\mathbb{E}} \mathbf{F}\) 有限，且此时 \(\dim \mathbf{F}' = \operatorname{codim}_{\mathbb{E}} \mathbf{F}\) .
\item
  \(\mathbf{F}'\) 在 \(\mathbf{E}\) 中的正交补等于 \(\mathbf{F}\) .
\item
  每个有限维子空间 \(\mathbf{G}'\) （ \(\mathbf{E}^*\) 的）都是 \(\mathbf{E}\) 的某个子空间的正交补，该子空间必然等于 \(\mathbf{G}'\) 在 \(\mathbf{E}\) 中的正交补，且具有有限余维数。
\item
  我们知道 \(\mathbf{F}'\) 同构于对偶 \((\mathbf{E} / \mathbf{F})^*\) （§ 2，第 6 号，命题 9 的推论），因此该断言由定理 4 得出，因为 \(\dim (\mathbf{E} / \mathbf{F}) = \operatorname{codim}_{\mathbf{E}}\mathbf{F}\) 根据定义。
\item
  设 \(F_{1}\) 是 \(F'\) 在 E 中的正交补；显然 \(F \subset F_{1}\) ，且 \(F_{1}'\) （ \(F_{1}\) 的正交补）等于 \(F'\) （§2，第4号）；典范线性映射 \((\mathrm{E}/\mathrm{F}_{1})^{*} \to (\mathrm{E}/\mathrm{F})^{*}\) （即 \(E/F \to E/F_{1}\) 的转置）因此是双射的（§2，第6号，命题10的推论）；于是由命题10可知，典范映射 \(E/F \to E/F_{1}\) 是双射的，这蕴含 \(F_{1} = F\) .
\item
  设 \(\mathbf{G}'\) 是 \(\mathbf{E}^*\) 的一个维数为有限数 \(p\) 的子空间，并设 \(\mathbf{F}\) 为它在 \(\mathbf{E}\) 中的正交补；则 \(\operatorname{codim}_{\mathbf{E}} \mathbf{F} \leqslant \dim \mathbf{G}'\) 。事实上，若 \((a_{i}^{*})_{1 \leqslant i \leqslant p}\) 是 \(\mathbf{G}'\) 的一组基，则 \(\mathbf{F}\) 是线性映射 \(x \mapsto (\langle x, a_{i}^{*} \rangle)\) （从 \(\mathbf{E}\) 到 \(\mathbf{K}_{s}^{p}\) ）的核，其秩至多为 \(p\) （第4号，命题9），由此得出结论（第4号）。然后设 \(\mathbf{F}'\) 是 \(\mathbf{F}\) 在 \(\mathbf{E}^*\) 中的正交补；由（i）可知 \(\dim \mathbf{F}' \leqslant \dim \mathbf{G}'\) ；但另一方面显然有 \(\mathbf{G}' \subset \mathbf{F}'\) ，因此 \(\mathbf{F}' = \mathbf{G}'\) ( \(\S 2\) ，第3号，命题4的推论4）。
\end{enumerate}

注记。一个无限维子空间 \(G'\) （含于 \(E^{*}\) ）不一定是 E 的某个子空间的正交补，换句话说，若 F 是 \(G'\) 在 E 中的正交补，则 F 的正交补 \(F'\) （在 \(E^{*}\) 中）可能不同于 \(G'\) （练习20(b)）。†

推论1. 设 \((x_{i}^{*})_{1\leqslant i\leqslant p}\) 是 E 上的线性型的一个有限序列，并设 F 为 E 中由使得下式成立的 \(x\) 所组成的子空间，

\[
\langle x, x _ {i} ^ {*} \rangle = 0 \quad for 1 \leqslant i \leqslant p.
\]

那么 \(\operatorname{codim}_{E} F\) 等于诸 \(x_{i}^{*}\) 之集的秩，并且 E 上每个在 F 上为零的线性型都是诸 \(x_{i}^{*}\) 的线性组合。进而 \(\operatorname{codim}_{E} F \leqslant p\) ，并且为使 \(\operatorname{codim}_{E} F = p\) ，必要且充分条件是诸 \(x_{i}^{*}\) 线性无关。

集合 \(G'\) （由诸 \(x_{i}^{*}\) 的线性组合组成）是 \(E^{*}\) 的子空间，且 F 是 \(G'\) 在 E 中的正交补，因此根据定理7， \(\operatorname{codim}_{E} F = \dim G'\) ；进一步 \(\dim G' \leqslant p\) ，而关系 \(\dim G' = p\) 意味着 \((x_{i}^{*})\) 是一个自由系（第2号，命题1）；由此得到该推论。

推论2. (i) 设 \((x_{i}^{*})_{1\leqslant i\leqslant p}\) 是 E 上的线性型的一个有限序列。要使 \((x_{i}^{*})\) 成为一个自由系，必要且充分条件是存在 E 的元素的一个序列 \((x_{i})_{1\leqslant i\leqslant p}\) 使得 \(\langle x_i,x_j^*\rangle = \delta_{ij}\) （克罗内克指标）。

\begin{enumerate}
\def\labelenumi{(\roman{enumi})}
\setcounter{enumi}{1}
\tightlist
\item
  设 \((x_{i})_{1\leqslant i\leqslant p}\) 是 E 的元素的一个有限序列。要使 \((x_{i})\) 成为一个自由系，必要且充分条件是存在 E 上的线性型的一个序列 \((x_{i}^{*})_{1\leqslant i\leqslant p}\) ，使得 \(\langle x_i,x_j^*\rangle = \delta_{ij}\) .
\end{enumerate}

显然，(ii) 由 (i) 得出，只需将 E 视为 \(\mathbf{E}^{**}\) 的一个子空间（借助 \(c_{\mathbf{E}}\) ，定理 6）。设 \(\mathbf{G}'\) 为 \(\mathbf{E}^*\) 中由诸 \(x_{i}^{*}\) 生成的子空间，而 \(\mathbf{F}\) 为它在 \(\mathbf{E}\) 中的正交补； \(\mathbf{E} / \mathbf{F}\) 和 \(\mathbf{G}'\) 中的每一个都可在典范意义下等同于另一个的对偶；若族 \((x_{i}^{*})\) 是自由的，则在 \(\mathbf{E} / \mathbf{F}\) 中存在一组基 \((\dot{x}_i)\) ，它是 \((x_{i}^{*})\) 的对偶基，而 \((x_{i})\) 作为诸类 \(\dot{x}_{i}\) 的任意代表系，都具有所需的性质。反之，系统 \((x_{i})\) （满足 \(\langle x_{i}, x_{j}^{*} \rangle = \delta_{ij}\) ）的存在意味着：对所有 \(i\) ， \(\mathbf{E}^*\) 中与 \(\mathbf{K}\) . \(x_{i}\) 正交的子空间包含诸 \(x_{j}^{*}\) （其指标 \(j \neq i\) ） ，但不包含 \(x_{i}^{*}\) ，因此系统 \((x_{i}^{*})_{1 \leqslant i \leqslant p}\) 是自由的。

推论3. 设 S 是一个集合，V 是由 S 到 K 的映射组成的右向量 K-空间 \(K_{d}^{S}\) 的一个向量子空间。要使 \(\dim V \geqslant p\) （其中 p 是整数），必要且充分条件是存在 S 的 p 个元素 \(s_{i}\) 以及 V 的 p 个元素 \(f_{i}\) （ \(1 \leqslant i \leqslant p\) ），使得 \(f_{i}(s_{j}) = \delta_{ij}\) .

该空间 \(K_{d}^{S}\) 典范地等同于 \(E = K_{d}^{(S)}\) 的对偶空间，且 \(f(s) = \langle e_{s}, f \rangle\) 对 \(s \in S\) 和 \(f \in K_{s}^{S}\) 成立 ，其中 \((e_{s})_{s \in S}\) 是 E 的典范基。推论2随后表明该条件是充分的。反之，假设 \(\dim V \geqslant p\) ，因此存在 V 的一个维数为 p 的子空间 \(G'\) ；设 F 为 \(G'\) 在 E 中的正交补，从而 \(\dim(E/F) = p\) 。由第 1 号定理 2 可知，存在 p 个元素 \(s_{i} \in S\) 使得 \(e_{s_{i}} (1 \leqslant i \leqslant p)\) 构成 F 在 E 中的一个补空间的一组基（将定理 2（第 1 号）应用于一个由 E 的一组基组成的自由子集，以及由此自由子集与 E 的典范基的并集组成的生成系）；然后我们取 \(f_{i}\) 为 G' 的一组基的元素，该基对偶于由诸 \(e_{s_{i}} \mod F\) 的类组成的 E/F 的基.

推论4. 设E是一个向量空间，M、N是E的两个有限余维子空间；若M'、N'分别是M和N在E*中的正交补，则M ∩ N在E*中的正交补为M' + N'。

由于 M（相应地，N）是 \(M'\) （相应地 \(N'\) ）在 E 中的正交补（定理 7）， \(M \cap N\) 是 \(M' + N'\) 在 E 中的正交补，因此 \(M' + N'\) 是 \(M \cap N\) 在 \(E^*\) 中的正交补（定理 7 (iii)）。

推论 5. 设 E 是维数为 n 的有限维向量空间。～对于 E 的每个维数为 p 的子空间 F，F 在 E* 中的正交补 F' 的维数为 n - p。～对于 E* 的每个维数为 q 的子空间 G'，G' 在 E 中的正交补 G 的维数为 n - q，且 G' 是 G 在 E* 中的正交补。

定理 7 给出了 E 中超平面的另一个刻画：

命题 11. 对于向量空间 E 中的每个超平面 H，存在 E 上的一个线性形式 \(x_0^*\) ，使得 \(\mathrm{H} = x_0^{*-1}(0)\) 。给定这样的一个形式 \(x_0^*\) ，对于 E 上的线性形式 \(x^*\) ，要满足 \(\mathrm{H} = x^{*-1}(0)\) ，必要且充分的条件是 \(x^* = x_0^*\alpha\) ，其中 \(\alpha\) 是一个标量 \(\neq 0\) 。反之，对于 E 上的每一个线性形式 \(x^* \neq 0\) ，子空间 \(x^{*-1}(0)\) 是 E 的一个超平面。

这个陈述仅仅表达了定理7对于E的余维数为1的子空间以及E*的维数为1的子空间的情形。

如果 H 是一个超平面，且 \(x_{0}^{*}\) 是满足 \(\mathrm{H}=x^{*-1}(0)\) 的一个线性型，则关系

\[
\langle x, x _ {0} ^ {*} \rangle = 0
\]

它刻画了元素 \(x \in \mathbf{H}\) ，称为 \(\mathbf{H}\) 的一个方程.

更一般地，如果 \((x_{\iota}^{*})\) 是 E 上的一族线性型，且 F 表示诸超平面 \(x_{\iota}^{*-1}(0)\) 的交所构成的向量子空间，则关系式“对所有 \(\iota\) , \(\langle x, x_{\iota}^{*}\rangle = 0\) 成立”刻画了 F 的元素；这些关系式

\[
\langle x, x _ {\iota} ^ {*} \rangle = 0 \quad \text { for   all } \iota
\]

构成子空间 F 的一个方程组。定理 7 (ii) 表达了 E 的每个向量子空间都可以由一个方程组来定义这一事实。

定理 7 (i) 和 (ii) 还证明了，余维数为有限数 p 的子空间 F 可以由 p 个方程组成的方程组

\[
\langle x, x _ {i} ^ {*} \rangle = 0, \quad 1 \leqslant i \leqslant p,\tag{19}
\]

来定义，其中这些型 \(x_{i}^{*}\) 是线性无关的。反之，定理 7 的推论 1 表明，子空间 \(F\) 若由 \(p\) 个方程 (19) 组成的方程组所定义，则其余维数 \(\leqslant p\) ，且其余维数为 \(p\) 当且仅当诸 \(x_{i}^{*}\) 线性无关；这等价于说 F 不能由至多包含 (19) 中 p - 1 个方程的方程组来定义。

\subsection{向量空间中的线性方程}

命题 12. 设 E, F 是域 K 上的两个向量空间，且 u:E → F 是一个线性映射。对于线性方程

\[
u (x) = y _ {0}\tag{20}
\]

至少有一个解 \(x \in \mathbf{E}\) ，必要且充分的条件是 \(y_0\) 与转置映射 \(^t u\) 的核正交.

\(u(\mathbf{E})\) 在 \(\mathbf{F}^*\) 中的正交补是 \({}^t u^{-1}(0)\) （§ 2，第 5 号，命题 8 的推论），而 \({}^t u^{-1}(0)\) 在 \(\mathbf{F}\) 中的正交补因此为 \(u(\mathbf{E})\) （第 5 号，定理 7 (ii)）。

我们将得到关于标量线性方程组

\[
\langle x, x _ {\iota} ^ {*} \rangle = \eta_ {\iota} \quad (\iota \in I)\tag{21}
\]

的一个更便捷的判别准则，其中未知量 x 在域 K 上的向量空间 E 中取值， \(x_{\iota}^{*}\) 是 E 上的线性型，而右端项 \(\eta_{\iota}\) 是 K 中的元素。如果我们考虑 E 的一组基

\((a_{\lambda})_{\lambda\in L}\) ，则方程组 (21) 等价于方程组

\[
\sum_ {\lambda \in L} \xi_ {\lambda} \left\langle a _ {\lambda}, x _ {\iota} ^ {*} \right\rangle = \eta_ {\iota} \quad (\iota \in I)\tag{22}
\]

其中 \(x = \sum_{\lambda \in L} \xi_{\lambda} a_{\lambda}\) ，而 (22) 的解必然是族 \((\xi_{\lambda})\) ，其元素属于 \(K\) 且具有有限支撑。

定义 4. \(\mathbf{E}^*\) 中由族 \((x_{\iota}^{*})\) 生成的子空间的维数称为系统 (21) 的秩。

命题 13. 系统 (21) 具有有限秩 \(r\) ，必要且充分条件是线性映射 \(u: x \mapsto (\langle x, x_{\iota}^* \rangle)\) 把 \(\mathbf{E}\) 映入 \(\mathbf{K}_s^I\) 且其秩为 \(r\) .

若 \(F'\) 是 \(E^{*}\) 中由诸 \(x_{\iota}^{*}\) 生成的子空间，则 u 的核是 \(F'\) 在 E 中的正交补 F；如果 \(F'\) 的维数为 r，则 F 的余维数为 r，反之亦然（第 5 号，定理 7），且 \(\operatorname{rg}(u) = \operatorname{codim}_{\mathbf{E}} \mathbf{F}\) （第 4 号，公式 (11)）。

定理 8. 设

\[
\langle x, x _ {\iota} ^ {*} \rangle = \eta_ {\iota} \quad (\iota \in I)\tag{21}
\]

设 E 是域 K 上的向量空间，考虑 E 上的一个标量线性方程组。要使该方程组至少有一个解，必须满足：对任意具有有限支撑的标量族 \((\rho_{\iota})\) ，只要 \(\sum_{\iota} x_{\iota}^{*}\rho_{\iota}=0\) ，就有 \(\sum_{\iota}\eta_{\iota}\rho_{\iota}=0\) 。如果系统 (21) 的秩是有限的，那么这个条件也是充分的。

该条件显然是必要的。它表明，如果 \(F'\) 是 \(E^{*}\) 中由族 \((x_{i}^{*})\) 生成的子空间，则存在一个线性映射 \(f:F'\to K_{d}\) ，使得 \(f(x_{\iota}^{*})=\eta_{\iota}\) 对所有 \(\iota\in I\) 成立。若 \(F'\) 是有限维的，维数为 r，则 \(F'\) 是 E 的余维数为 r 的子空间 F 的正交补（第 5 号，定理 7），并且 \(F'\) 等同于 E/F 的对偶空间（§ 2，第 6 号，命题 9 的推论）；因此 f 是双对偶空间 \((\mathrm{E}/\mathrm{F})^{**}\) 中的一个元素。由于 E/F 是有限维的，存在唯一一个元素 \(y\in E/F\) ，使得 \(f(x^{*})=\langle y,x^{*}\rangle\) 对所有 \(x^{*}\in F'\) 成立（第 5 号，定理 6）。于是 (21) 的解即为那些在 E/F 中的典范像为 y 的 \(x\in E\) 。

注记。当系统 (21) 的秩为无穷时，定理 8 的条件不再充分。例如，假设 \(x_{\iota}^{*}\) 是空间 \(\mathbf{E} = \mathbf{K}_{s}^{(\mathrm{I})}\) 上的坐标线性型，其中 I 是无限的（§ 2，第 6 号）；由于 \(x_{\iota}^{*}\) 线性无关，定理 8 的条件对每个族 ( \(\eta_{\iota}\) ) 都成立，但系统 (21) 仅当族 ( \(\eta_{\iota}\) ) 具有有限支撑时才有解。

系统 (21) 若只含有限个方程，则其秩总是有限的，且此时秩至多等于方程的个数（第 2 号，命题 1）。类似地，若 E 为有限维，维数为 n（对于系统 (22) 而言，这对应于未知量仅有有限个 n 的情形），则其对偶 \(E^{*}\) 的维数为 n，因此系统 (21) 的秩至多为 n（第 3 号，命题 4 的推论 4）。由此我们推出：

推论 1. 向量空间中由有限个方程组成的标量线性方程组，若各方程的左端是线性无关的线性型，则恒有解。

推论 2. 对于系数在域 K 中、含 n 个未知量的齐次线性方程组 (22)，要使它存在由 K 中元素组成的非平凡解，必要且充分条件是其秩 \textless n.

如果方程的个数是有限的并且 \textless n ，那么情况总是如此。

推论 3. 对于系数和右端项均属于域 K 的线性方程组 (22)，若其包含 n 个方程和 n 个未知数，则它有且仅有一个由 K 中元素构成的解，当且仅当相应的齐次方程组没有非平凡解（或者，等价地说，该方程组的左端是线性无关的线性型）。

\subsection{向量空间的张量积}

第3、4、5节中关于自由模或投射模的结果特别适用于向量空间，并给出以下性质：

命题14. 给定一个正合序列

\[
0 \rightarrow \mathbf {E} ^ {\prime} \rightarrow \mathbf {E} \rightarrow \mathbf {E} ^ {\prime \prime} \rightarrow 0\tag{23}
\]

关于域 K 上的右向量空间与线性映射，以及域 K 上的左向量空间 F，相应的 Z-线性映射序列

\[
0 \rightarrow \mathrm{E} ^ {\prime} \otimes_ {\mathbb {K}} \mathrm{F} \rightarrow \mathrm{E} \otimes_ {\mathbb {K}} \mathrm{F} \rightarrow \mathrm{E} ^ {\prime \prime} \otimes_ {\mathbb {K}} \mathrm{F} \rightarrow 0
\]

是正合的且分裂。

由于序列 (23) 分裂，这是 §3 第 7 号命题 7 的推论 5 以及 §3 第 6 号命题 5 的一个特例。

由于命题 14，当 \(\mathbf{E}'\) 是 \(\mathbf{E}\) 的一个向量子空间， \(j: \mathbf{E}' \to \mathbf{E}\) 为典范单射时， \(\mathbf{E}' \otimes_{\mathbb{K}} \mathbf{F}\) 通常作为 \(\mathbf{Z}\) -模被等同于 \(\mathbf{E} \otimes_{\mathbb{K}} \mathbf{F}\) 的一个子模，所借助的是单射 \(j \otimes 1_{\mathbf{F}}\) 。在此约定下：

推论。设 K 为域，E 为 K 上的右向量空间，F 为 K 上的左向量空间， \((\mathbf{M}_{\alpha})_{\alpha \in \mathbf{A}}\) 为 E 的一族向量子空间， \((\mathbf{N}_{\beta})_{\beta \in \mathbf{B}}\) 为 F 的一族向量子空间。则

\[
\left(\bigcap_ {\alpha \in \mathrm{A}} \mathrm{M} _ {\alpha}\right) \otimes_ {\mathrm{K}} \left(\bigcap_ {\beta \in \mathrm{B}} \mathrm{N} _ {\beta}\right) = \bigcap_ {(\alpha , \beta) \in \mathrm{A} \times \mathrm{B}} \left(\mathrm{M} _ {\alpha} \otimes_ {\mathrm{K}} \mathrm{N} _ {\beta}\right).\tag{24}
\]

显然只需证明该特例

\[
\left(\bigcap_ {\alpha \in \mathrm{A}} \mathrm{M} _ {\alpha}\right) \otimes_ {\mathrm{K}} \mathrm{F} = \bigcap_ {\alpha \in \mathrm{A}} \left(\mathrm{M} _ {\alpha} \otimes_ {\mathrm{K}} \mathrm{F}\right).\tag{25}
\]

显然，(25) 的左端包含于右端。为证明反向包含关系，我们考虑 F 的一组基 \((f_{\lambda})_{\lambda \in L}\) 。 \(\mathbf{E} \otimes_{\mathbb{K}} \mathbf{F}\) 的每个元素于是都可以唯一地表示为 \(\sum_{\lambda \in L} x_{\lambda} \otimes f_{\lambda}\) 的形式，其中 \(x_{\lambda} \in \mathbf{E}\) （§ 3，第 7 号，命题 7 的推论 1）；如果 \(\mathbf{E}'\) 是 E 的一个向量子空间，则关系式 \(\sum_{\lambda \in L} x_{\lambda} \otimes f_{\lambda} \in \mathbf{E}' \otimes_{\mathbb{K}} \mathbf{F}\) 依命题 14 等价于 \(x_{\lambda} \in \mathbf{E}'\) 对所有 \(\lambda \in L\) 成立。说 \(\sum_{\lambda \in L} x_{\lambda} \otimes f_{\lambda}\) 属于每一个 \(\mathbf{M}_{\alpha} \otimes_{\mathbb{K}} \mathbf{F}\) ，这就意味着对所有 \(\lambda \in L\) 以及所有 \(\alpha \in A\) ， \(x_{\lambda} \in M_{\alpha}\) ，即 \(x_{\lambda} \in \bigcap_{\alpha \in A} M_{\alpha}\) 对所有 \(\lambda \in L\) 成立，这证明了 (25) 的右端包含于左端。命题 15. 若 \((E_{\lambda})_{\lambda \in L}\) 是域 K 上的一族右向量空间，且 \((F_{\mu})_{\mu \in M}\) 是 K 上的一族左向量空间，则典范映射

\[
\left(\prod_ {\lambda \in L} E _ {\lambda}\right) \otimes_ {K} \left(\prod_ {\mu \in M} F _ {\mu}\right)\rightarrow \prod_ {(\lambda , \mu) \in L \times M} (E _ {\lambda} \otimes_ {K} F _ {\mu})\tag{26}
\]

(§ 3，第 7 号，公式 (22)) 是单射的。

我们写 \(\mathbf{F} = \prod_{\mu \in \mathbf{M}} \mathbf{F}_{\mu}\) ；映射 (26) 是诸典范映射的复合。

\[
\left(\prod_ {\lambda \in L} \mathbf {E} _ {\lambda}\right) \otimes_ {\mathbb {K}} F \rightarrow \prod_ {\lambda \in L} \left(\mathbf {E} _ {\lambda} \otimes_ {\mathbb {K}} F\right)
\]

和

\[
\prod_ {\lambda \in L} \left(\mathrm{E} _ {\lambda} \otimes_ {K} F\right)\rightarrow \prod_ {\lambda \in L} \left(\prod_ {\mu \in M} \left(\mathrm{E} _ {\lambda} \otimes_ {K} F _ {\mu}\right)\right);
\]

由于 F 和诸 \(\mathbf{E}_{\lambda}\) 是 K 上的向量空间，这就归结为 § 3，第 7 号，命题 7 的推论 3。

当命题 15 的条件满足时，张量积 \(\left(\prod_{\lambda \in L} E_{\lambda}\right) \otimes_{K} \left(\prod_{\mu \in M} F_{\mu}\right)\) 通常等同于它在以下空间中的典范像

\[
\prod_ {\lambda , \mu} (\mathbf {E} _ {\lambda} \otimes_ {K} \mathbf {F} _ {\mu}).
\]

按照这个约定：

推论。设 F 是 K 上的左向量空间；对于每个集合 X，左向量空间 \(K_{d}^{X} \otimes_{K} F\) 等同于空间 \(F^{X}\) （即由 X 到 F 的所有映射组成的空间）的子空间，它由满足以下条件的映射 u 组成： \(u(X)\) 在 F 中具有有限秩。

若 \((f_{\lambda})\) 是 \(F\) 的一组基，则 \(\sum_{\lambda \in L} v_{\lambda} \otimes f_{\lambda}\) 作为 \(K_d^X \otimes_K F\) 的元素在 (26) 下等同于映射 \(x \mapsto \sum_{\lambda} v_{\lambda}(x)f_{\lambda}\) 。由于 \(v_{\lambda} = 0\) ，除了指标 \(\lambda\) 属于一个有限子集 \(H\) （该子集含于 \(L\) ）以外， \(X\) 在上述映射下的像包含于 \(F\) 中由诸 \(f_{\lambda}\) （其指标 \(\lambda \in H\) ）生成的有限维子空间。反之，设 \(u: X \to F\) 为一个映射，使得 \(u(X)\) 包含于一个有限维子空间 \(G\) （该子空间含于 \(F\) ），并设 \((b_i)_{1 \leqslant i \leqslant n}\) 是 \(G\) 的一组基。对所有 \(x \in X\) ，我们可以写成 \(u(x) = \sum_{i=1}^{n} v_i(x)b_i\) ，其中诸 \(v_i(x)\) 是 \(K\) 的完全确定的元素；因此 \(n\) 个映射 \(v_i: X \to K\) 被定义，且显然 \(u\) 随后与元素 \(\sum_{i=1}^{n} v_i \otimes b_i\) 等同.

类似地，对于域 K 上的右向量空间 E 和集合 Y， \(E \otimes_{K} K_{s}^{Y}\) 等同于空间 \(E^{Y}\) 的一个子空间，它由满足以下条件的映射 \(v: Y \to E\) 组成： \(v(Y)\) 具有有限秩。更特别地，对于每个域 K， \(K_{d}^{X} \otimes_{K} K_{s}^{Y}\) 等同于空间 \(K^{X \times Y}\) （由 \(X \times Y\) 到 K 中的映射组成）的一个子空间（K 被视为 (K, K)-双模）；元素 \(\sum_{i} u_{i} \otimes v_{i}\) （其中 \(u_{i}\) 是 X 到 K 的映射，而 \(v_{i}\) 是 Y 到 K 的映射）被等同于映射 \((x, y) \mapsto \sum_{i} u_{i}(x) v_{i}(y)\) ，它是 \(X \times Y\) 到 K 中的映射。

命题 16. (i) 设 K, L 为两个域，E 为 K 上的左向量空间，F 为 L 上的左向量空间，G 为 (K, L)-双模。则典范 Z-同态

\[
\nu : \operatorname{Hom} _ {\mathbf {K}} (\mathrm{E}, \mathrm{G}) \otimes_ {\mathrm{L}} \mathrm{F} \rightarrow \operatorname{Hom} _ {\mathbf {K}} (\mathrm{E}, \mathrm{G} \otimes_ {\mathrm{L}} \mathrm{F})\tag{27}
\]

（§ 4，第 2 号，公式 (7)）是单射；当向量空间 E, F 中有一个是有限维时，它是双射。

\begin{enumerate}
\def\labelenumi{(\roman{enumi})}
\setcounter{enumi}{1}
\tightlist
\item
  设 \(E_{1}\) , \(E_{2}\) , \(F_{1}\) , \(F_{2}\) 为交换域 K 上的四个向量空间；则典范 K-同态
\end{enumerate}

\[
\lambda : \operatorname{Hom} \left(\mathrm{E} _ {1}, \mathrm{F} _ {1}\right) \otimes \operatorname{Hom} \left(\mathrm{E} _ {2}, \mathrm{F} _ {2}\right)\rightarrow \operatorname{Hom} \left(\mathrm{E} _ {1} \otimes \mathrm{E} _ {2}, \mathrm{F} _ {1} \otimes \mathrm{F} _ {2}\right)\tag{28}
\]

(§ 4，第 4 号，公式 (21)) 是单射；若有序对 \((\mathbf{E}_1, \mathbf{E}_2), (\mathbf{E}_1, \mathbf{F}_1), (\mathbf{E}_2, \mathbf{F}_2)\) 之一由有限维空间组成，则它是双射。

断言 (i) 是 § 4 第 2 号命题 2 的特殊情形。类似地，(ii) 的第二个断言是 § 4 第 4 号命题 4 的特殊情形。最后，要看出同态 (28) 总是单射的，注意 \(\operatorname{Hom}(\mathbf{E}_i, \mathbf{F}_i)\) 是 \(\mathbf{F}_i^{\mathbf{E}_i}\) 的一个向量子空间 ( \(i = 1, 2\) )，并且

\[
\operatorname{Hom} \left(\mathrm{E} _ {1} \otimes \mathrm{E} _ {2}, \mathrm{F} _ {1} \otimes \mathrm{F} _ {2}\right)
\]

典范地等同于空间 \((\mathbf{F}_{1} \otimes \mathbf{F}_{2})^{\mathbf{E}_{1} \times \mathbf{E}_{2}}\) 的一个向量子空间（II，§ 3，第 1 号，命题 1）；当完成这些等同，并且 (28) 的左端也被等同于 \(F_{1}^{E_{1}} \otimes F_{2}^{E_{2}}\) 的一个子空间（命题 14）时，典范映射 (28) 就变为典范映射 (26) 在该子空间上的限制，并且已经看到（命题 15）后者是单射的。

推论 1. 设 E 和 F 是域 K 上的两个向量空间；则典范映射

\[
\mathbf {E} ^ {*} \otimes_ {\mathbb {K}} \mathbf {F} \rightarrow \operatorname{Hom} _ {\mathbb {K}} (\mathbf {E}, \mathbf {F})
\]

(§ 4，第 2 号，公式 (11)) 是单射；当 E 或 F 为有限维时，它是双射。

这是命题 16 (i) 的一个特例。

推论 2. 设 E 是右向量空间，F 是同一域 K 上的左向量空间；则典范映射

\[
\mathbf {E} \otimes_ {\mathbb {K}} \mathbf {F} \rightarrow \operatorname{Hom} _ {\mathbb {K}} (\mathbf {E} ^ {*}, \mathbf {F})\tag{29}
\]

（§ 4，第 2 号，公式 (15)）是单射；当 E 是有限维时，它是双射。

这是 § 4 第 2 号注记 2 的一个特例。

注记。(1) 设 K 是一个交换域，E、F 是 K 上的两个向量空间，

\((a_{\lambda})\) 是 \(\mathbf{E}\) 的一组基， \((b_{\mu})\) 是 \(\mathrm{F}\) 的一组基；于是 \((a_{\lambda} \otimes b_{\mu})\) 是向量 \(\mathbf{K}\) -空间 \(\mathbf{E} \otimes_{\mathbb{K}} \mathbf{F}\) 的一组基（II，§ 3，第 7 号，命题 2 的推论 2），因此

\[
\dim_ {\mathbf {K}} (\mathbf {E} \otimes_ {\mathbf {K}} \mathbf {F}) = \dim_ {\mathbf {K}} \mathbf {E}. \dim_ {\mathbf {K}} \mathbf {F}.\tag{30}
\]

\begin{enumerate}
\def\labelenumi{(\arabic{enumi})}
\setcounter{enumi}{1}
\tightlist
\item
  设 K 是一个交换域， \(E_{1}\) , \(E_{2}\) , \(F_{1}\) , \(F_{2}\) 为 K 上的四个向量空间， \(u: E_{1} \rightarrow F_{1}\) , \(v: E_{2} \rightarrow F_{2}\) 为两个线性映射；则
\end{enumerate}

\[
\operatorname{rg} (u \otimes v) = \operatorname{rg} (u). \operatorname{rg} (v).\tag{31}
\]

显然 \((u\otimes v)(\mathbf{E}_1\otimes \mathbf{E}_2)\) 是 \(u(\mathbf{E}_1)\otimes v(\mathbf{E}_2)\) 在 \(\mathbf{F}_1\otimes \mathbf{F}_2\) 中的典范像，因此（命题 14）同构于 \(u(\mathbf{E}_1)\otimes v(\mathbf{E}_2)\) ；结论随后由 (30) 得出。

\begin{enumerate}
\def\labelenumi{(\arabic{enumi})}
\setcounter{enumi}{2}
\tightlist
\item
  在注记 1 的相同假设下，
\end{enumerate}

\[
\dim_ {\mathbb {K}} (\operatorname{Hom} _ {\mathbb {K}} (\mathrm{E}, \mathrm{F})) \geqslant \dim_ {\mathbb {K}} \mathrm{E}. \dim_ {\mathbb {K}} \mathrm{F}.\tag{32}
\]

如果 E 同构于 \(K^{(I)}\) ，则 \(\operatorname{Hom}(E, F)\) 同构于 \((\operatorname{Hom}(K, F))^{I}\) （§ 1，第 6 号，命题 6 的推论 1），从而同构于 \(F^{I}\) （§ 1，第 14 号）；由于 \(F^{(I)}\) 是 \(F^{I}\) 的一个子空间，且 \(\dim(F^{(I)}) = \operatorname{Card}(I).dim F = \dim E. dim F\) （第 2 号，命题 2），不等式 (32) 由第 3 号命题 4 的推论 4 得出。同样的论证表明，当 dim E 有限时，(32) 的两边相等（参见 § 10，第 3 和第 4 号）。

\subsection{张量积中元素的秩}

设 E 是右向量空间，F 是同一域 K 上的左向量空间；对每个元素 \(u \in E \otimes_{K} F\) ，在 (29) 下典范地对应一个同态 \(u_{1} \in \operatorname{Hom}_{\mathbb{K}}(\mathbf{E}^{*}, \mathbf{F})\) ；若 \(u = \sum_{i} x_{i} \otimes y_{i}\) ，其中 \(x_{i} \in E, y_{i} \in F\) ，则元素 \(u_{1}\) 是线性映射

\[
x ^ {*} \mapsto \sum_ {i} \langle x ^ {*}, x _ {i} \rangle y _ {i}.\tag{33}
\]

另一方面， \(\mathbf{E} \otimes_{\mathbb{K}} \mathbf{F}\) 被典范地等同于 \(\mathbf{F} \otimes_{\mathbb{K}^0} \mathbf{E}\) ，其中 \(\mathbf{E}\) 被视为左向量空间，而 \(\mathbf{F}\) 被视为反域 \(\mathbf{K}^0\) 上的右向量空间；因此典范地对应于 \(u\) 一个同态 \(u_2 \in \operatorname{Hom}_{\mathbb{K}}(\mathbf{F}^*, \mathbf{E})\) ，它由下式给出

\[
y ^ {*} \mapsto \sum_ {i} x _ {i} \langle y _ {i}, y ^ {*} \rangle ;\tag{34}
\]

\(u_{1}\) （相应地 \(u_{2}\) ）被视为 \(E^{*}\) 到 \(F^{**}\) 中（相应地， \(F^{*}\) 到 \(E^{**}\) 中）的映射时，正是 \(u_{2}\) （相应地 \(u_{1}\) ）的转置。因此 \(u_{1}\) 和 \(u_{2}\) 的秩等于同一个有限数 r，即 F 的子空间 \(u_{1}(E^{*})\) 与 E 的子空间 \(u_{2}(F^{*})\) 的公共维数，其中每一个都典范同构于另一个的对偶（第 5 号）；我们称 r（记为 \(\operatorname{rg}(u)\) ）为 \(E \otimes_{K} F\) 的元素 u 的秩，并称 \(u_{1}(E^{*})\) 和 \(u_{2}(F^{*})\) 为与 u 相关联的子空间（分别属于 F 和 E）。

命题 17. 设 \(u\) 是 \(\mathbf{E} \otimes_{\mathbb{K}} \mathbf{F}\) 的一个元素，而 \(\mathbf{M} \subset \mathbf{E}\) 和 \(\mathbf{N} \subset \mathbf{F}\) 是与其相关联的子空间。对每一个表达式 \(u = \sum_{i=1}^{s} x_i \otimes y_i\) （ \(u\) 的表达式），其中 \(x_i \in \mathbf{E}\) 和 \(y_i \in \mathbf{F}\) （ \(1 \leqslant i \leqslant s\) ），子空间 \(\mathbf{M}\) （相应地 \(\mathbf{N}\) ) 包含在 \(\mathbf{E}\) （相应地 \(\mathbf{F}\) ) 的由诸 \(x_i\) （相应地诸 \(y_i\) ) 生成的子空间中。此外，以下性质是等价的：

\begin{enumerate}
\def\labelenumi{(\alph{enumi})}
\item
  整数 \(s\) 等于 \(u\) 的秩.
\item
  族 \((x_{i})_{1\leqslant i\leqslant s}\) 是M的一组基。
\item
  族 \((y_{i})_{1\leqslant i\leqslant s}\) 是 N 的一组基。
\item
  族 \((x_{i})_{1\leqslant i\leqslant s}\) 和 \((y_{i})_{1\leqslant i\leqslant s}\) 都是自由的。
\end{enumerate}

由 (33)（相应地 (34)）， \(N = u_{1}(E^{*})\) （相应地 \(M = u_{2}(F^{*})\) ) 的每个元素都是诸 \(y_{i}\) （相应地诸 \(x_{i}\) ）的线性组合；由此得第一个断言。若 s = r，则由诸 \(x_{i}\) （相应地 \(y_{i}\) ）生成的、维数 \(\leqslant \dim M\) （相应地 \(\leqslant \dim N\) ）且包含 M（相应地 N）的子空间与它相同，因此 (a) 蕴含 (b) 和 (c)，更不必说 (d)。反之，条件 (b) 和 (c) 中的每一个都依 \(\operatorname{rg}(u)\) 的定义蕴含 (a)。最后，若 (d) 成立，则存在一个族 \((x_{i}^{*})_{1\leqslant i\leqslant s}\) ，其元素属于 \(E^{*}\) ，使得 \(\langle x_{i}, x_{j}^{*}\rangle = \delta_{ij}\) （第 5 号，定理 7 的推论 1），因此由 (33) 可得 \((y_{i})\) 是 N 的一组基，这就完成了证明。

推论 1. \(u\) 的秩是最小整数 \(s\) ，使得存在一个表达式 \(u = \sum_{i=1}^{s} x_i \otimes y_i\) ，其中 \(x_i \in \mathbf{E}\) 和 \(y_i \in \mathbf{F}\) 对于 \(1 \leqslant i \leqslant s\) .

这直接由命题 17 及第 2 号命题 1 得出。

推论 2. 设 K 为一个交换域，E、F 为 K 上的两个向量空间，L 为 K 的一个交换扩域。设 u 为 E ⊗K F 的一个元素，M 和 N 为其相伴子空间，u' 为 u 在 (E ⊗K F)(L) 中的典范像（典范地等同于 E(L) ⊗L F(L)，参见 § 5 第 1 号命题 3）；则 rg(u') = rg(u)，且与 u' 相伴的子空间典范地等同于 M(L) 和 N(L)。

若 \(u = \sum_{i=1}^{r} x_i \otimes y_i\) ，其中族 \((x_i)\) 和 \((y_i)\) 是自由的，则 \(u' = \sum_{i=1}^{r} (1 \otimes x_i) \otimes (1 \otimes y_i)\) 且族 \((1 \otimes x_i)\) 和 \((1 \otimes y_i)\) 分别在 \(E_{(L)}\) 和 \(F_{(L)}\) 中是自由的（§ 5，第 1 号，命题 4）。

\subsection{向量空间的标量扩张}

回顾（I，§ 9，第 1 号，定理 2），域 K 到非零环 A 的同态必然是单射。

命题 18. 设 \(\rho\) 为域 \(K\) 到环 \(A\) 中的一个同态。对于每一个由向量 \(K\) -空间和 \(K\) -线性映射组成的正合序列

\[
\mathrm{E} ^ {\prime} \xrightarrow {u} \mathrm{E} \xrightarrow {v} \mathrm{E} ^ {\prime \prime}
\]

该序列

\[
\mathrm{E} _ {(\mathrm{A})} ^ {\prime} \xrightarrow {u _ {(\mathrm{A})}} \mathrm{E} _ {(\mathrm{A})} \xrightarrow {v _ {(\mathrm{A})}} \mathrm{E} _ {(\mathrm{A})} ^ {\prime \prime}
\]

是正合的。

这是第 7 号命题 14 的一个特例，需考虑到 § 1 第 4 号中的注记 4。

推论。对于每一个K-线性映射 \(f: \mathrm{E}' \to \mathrm{E}\) , \(\operatorname{Im}(f_{(\mathbf{A})}) = (\operatorname{Im}(f))_{(\mathbf{A})}\) , \(\operatorname{Ker}(f_{(\mathbf{A})}) = (\operatorname{Ker}(f))_{(\mathbf{A})}\) , \(\operatorname{Coker}(f_{(\mathbf{A})}) = (\operatorname{Coker}(f))_{(\mathbf{A})}\) ，在典范同构的意义下。

命题 19. 设 \(\rho\) 为域 K 到环 A 中的一个单同态。对于 K 上的每个左向量空间 E，典范映射 \(\phi: \mathbf{E} \to \rho^{*}(\mathbf{E}) = \mathbf{A} \otimes_{\mathbb{K}} \mathbf{E}\) 是单射。此外，对 E 的每个向量子空间 \(\mathbf{E}'\) ， \(\rho^{*}(\mathbf{E}') = \mathbf{A} \otimes_{\mathbb{K}} \mathbf{E}'\) 典范地等同于 \(\mathbf{A} \otimes_{\mathbb{K}} \mathbf{E}\) 的一个直因子子 A-模，并且在这种等同下，

\[
(\mathbf {A} \otimes_ {\mathbf {K}} \mathbf {E} ^ {\prime}) \cap \phi (\mathbf {E}) = \phi (\mathbf {E} ^ {\prime}).\tag{35}
\]

第一个断言是 § 5 第 1 号命题 4 的特例；第二个是第 7 号命题 14 的特例；最后，为证明 (35)，只需在 A（视为右 K-模）中取一组基 \((a_{\lambda})_{\lambda \in L}\) ，使得 \(a_{\lambda_0} = 1\) 对某个指标 \(\lambda_0\) 成立（第 1 号，定理 2）；A \(\otimes_{\mathbb{K}}\) E 的元素可以唯一地写成 \(\sum_{\lambda} a_{\lambda} \otimes x_{\lambda}\) ，其中 \(x_{\lambda} \in \mathbf{E}\) ，而这样的元素属于 A \(\otimes_{\mathbb{K}}\) E' 的必要且充分条件是： \(x_{\lambda} \in \mathbf{E}'\) 对所有 \(\lambda\) 成立。另一方面， \(\phi(\mathbf{E})\) 的元素是那些使得 \(x_{\lambda} = 0\) 对 \(\lambda \neq \lambda_0\) 成立的元素；对元素 \(\sum_{\lambda} a_{\lambda} \otimes x_{\lambda}\) 属于 (A \(\otimes_{\mathbb{K}}\) E') \(\cap\) \(\phi(\mathbf{E})\) ，必要且充分的条件是 \(x_{\lambda} = 0\) 对 \(\lambda \neq \lambda_0\) 成立，且 \(x_{\lambda_0} \in \mathbf{E}'\) ，由此得出结论。

推论。设 \(\rho\) 为域 K 到环 A 中的一个单同态。K-线性映射 \(f: \mathbf{E} \to \mathbf{F}\) （其中 E 和 F 是 K 上的两个向量空间）为单射（相应地为满射、零映射），当且仅当 \(f_{(\mathbf{A})}: \mathbf{E}_{(\mathbf{A})} \to \mathbf{F}_{(\mathbf{A})}\) 为单射（相应地为满射、零映射）。

这直接由命题 19 及命题 18 的推论得出。

命题 20. 设 \(\rho\) 为域 \(\mathbf{K}\) 到环 \(\mathbf{A}\) 中的一个同态。对于每一个左向量空间 \(\mathbf{E}\) （ \(\mathbf{K}\) 上的），典范右 A-模同态

\[
\upsilon : \left(\mathrm{E} ^ {*}\right) _ {(\mathrm{A})} \rightarrow \left(\mathrm{E} _ {(\mathrm{A})}\right) ^ {*}
\]

（§ 5，第 4 号）是单射；当 E 为有限维时，它是双射。

第二个断言由 § 5 第 4 号命题 8 得出。为证明第一个断言，我们注意到 \((\mathbf{E}^{*})_{(\mathbf{A})}\) 的每个元素都可以写成 \(\sum_{i} x_{i}^{*} \otimes \alpha_{i}\) ，其中 \(\alpha_{i} \in \mathbf{A}\) ，而 \((x_{i}^{*})_{1 \leqslant i \leqslant n}\) 是 \(\mathbf{E}^{*}\) 中的一个自由族；在 \((\mathbf{E}_{(\mathbf{A})})^{*}\) 中对应于它的是线性形式 \(y^{*}\) ，使得 \(y^{*}(1 \otimes x) = \sum_{i} \rho(\langle x, x_{i}^{*} \rangle)\alpha_{i}\) 对所有 \(x \in \mathbf{E}\) 成立。但是，E 中存在一个族 \((x_{i})_{1 \leqslant i \leqslant n}\) ，使得 \(\langle x_{i}, x_{j}^{*} \rangle = \delta_{ij}\) （第 5 号，定理 7 的推论 2），由此 \(y^{*}(1 \otimes x_{i}) = \alpha_{i}\) ；因此关系 \(y^{*} = 0\) 蕴含 \(\alpha_{i} = 0\) 对所有 \(i\) 成立，这证明了我们的断言。

命题 21. 设 K 是一个域，L 是 K 的一个扩域。

\begin{enumerate}
\def\labelenumi{(\roman{enumi})}
\item
  对于每个向量空间 \(\mathbf{E}\) （ \(\mathbf{K}\) 上的）， \(\dim_{\mathbf{L}}(\mathrm{E}_{(\mathbf{L})}) = \dim_{\mathbf{K}}\mathbf{E}\) .
\item
  对于每个 K-线性映射 \(u: \mathbf{E} \to \mathbf{F}\) （其中 \(\mathbf{E}\) 和 \(\mathbf{F}\) 是 \(\mathbf{K}\) 上的向量空间）， \(\operatorname{rg}(u_{(L)}) = \operatorname{rg}(u)\) .
\end{enumerate}

若 \((e_{\iota})_{\iota\in I}\) 是 \(E\) 在 \(K\) 上的一组基，则 \((1\otimes e_{\iota})_{\iota\in I}\) 是 \(\mathbf{E}_{(L)}\) 在 \(L\) 上的一组基（§ 5，第 1 号，命题 4），由此得第一个断言；第二个断言由第一个断言及以下事实推出： \(u_{(L)}(\mathbf{E}_{(L)})\) 依命题 18 的推论典范地等同于 \((u(\mathbf{E}))_{(L)}\) 。

命题 22. 设 K 是一个交换域， \(\rho: \mathrm{K} \to \mathrm{A}\) 为一个单射的中心同态，且 E、F 是 K 上的两个向量空间。则典范同态

\[
\omega : \mathrm{A} \otimes_ {\mathrm{K}} \operatorname{Hom} (\mathrm{E}, \mathrm{F}) \rightarrow \operatorname{Hom} _ {\mathrm{A}} \left(\mathrm{E} _ {(\mathrm{A})}, \mathrm{F} _ {(\mathrm{A})}\right)\tag{36}
\]

（§ 5，第 3 号，公式 (17)）是单射的；如果 A 或 E 是 K 上的有限维向量空间，则它是双射的。

这是 § 5 第 3 号命题 7 的一个特例。

\subsection{整环上的模}

命题23. 在整环 A 上的模 E 中，非自由元素的集合 T 是 E 的一个子模。

若 x 与 y 均非自由，则存在 A 中两个非零元素 \(\alpha\) , \(\beta\) ，使得 \(\alpha x = 0\) 和 \(\beta y = 0\) 。则由于 A 是整环， \(\alpha\beta \neq 0\) ，且由于 A 是交换的， \(\alpha\beta(\lambda x + \mu y) = 0\) 对 A 中所有 \(\lambda\) 和 \(\mu\) 成立 ，因此 \(\lambda x + \mu y\) 不是自由的。

注记。设 E 是任意交换环 A 上的模。若 x 是 E 的非自由元素，则子模 Ax 中的每个元素都是非自由的。另一方面，若 A 含有零因子，则 E 中两个非自由元素之和可能是自由的；例如，在 Z/6Z（视为自身上的模）中，3 和 4 不是自由的，但 \(3 + 4 = 1\) 是自由的。

命题23引出以下定义：

定义5. 在整环A上的模E中，E的挠子模是由非自由元素（也称为E的挠元素）组成的E的子模。

当 E 等于其挠子模时（即 E 的每个元素都被 A 的一个元素 \(\neq0\) 所零化时），E 被称为挠模。当 E 的挠子模约化为 0（即 E 的每个非零元素都是自由的）时，E 被称为（按语言习惯的滥用）无挠模。

自由 A-模的每个子模（特别地，每个投射 A-模）都是无挠的。Z-模 Q 是无挠的。

命题 24. 设 A 是一个整环。对于每个 A-模 E，令 T(E) 表示 E 的挠子模。设 \(f: \mathbf{E} \to \mathbf{E}'\) 是一个 A-线性映射，E 和 \(\mathbf{E}'\) 都是 A-模。

\begin{enumerate}
\def\labelenumi{(\roman{enumi})}
\item
  \(f(\mathbf{T}(\mathbf{E})) \subset \mathbf{T}(\mathbf{E}')\) .
\item
  若 \(f\) 是单射，则 \(f(\mathbf{T}(\mathbf{E})) = \mathbf{T}(\mathbf{E}') \cap f(\mathbf{E})\) .
\item
  若 \(f\) 是满射且 \(\operatorname{Ker}(f) \subset \mathrm{T}(\mathbf{E})\) ，则 \(f(\mathrm{T}(\mathbf{E})) = \mathrm{T}(\mathbf{E}')\) .
\end{enumerate}

断言 (i) 和 (ii) 是显然的。另一方面，如果 \(f\) 是满射且 \(x' \in \mathrm{T}(\mathrm{E}')\) ，则 \(x' = f(x)\) ，其中 \(x \in \mathbf{E}\) ，且根据假设存在 \(\alpha \neq 0\) 于 \(\mathbf{A}\) 中，使得 \(f(\alpha x) = \alpha x' = 0\) ；由此 \(\alpha x \in \operatorname{Ker}(f)\) ，且由假设存在 \(\beta \neq 0\) 于 \(\mathbf{A}\) 中，使得 \(\beta(\alpha x) = 0\) ；由于 \(\beta \alpha \neq 0\) ，故 \(x \in \mathrm{T}(\mathrm{E})\) .

推论 1. 对于每个 A-模 E，E/T(E) 是无挠的。

若 \(f: \mathbf{E} \to \mathbf{E}'\) 是一个 A-线性映射，令 \(f_{\mathrm{T}}\) 表示映射 \(\mathrm{T}(\mathbf{E}) \to \mathrm{T}(\mathbf{E}')\) ，它的图与 \(f\) 在 \(\mathrm{T}(\mathbf{E})\) 上的限制的图相同。利用此记号：

推论 2. 对于 A-模与 A-线性映射的每个正合序列

\[
0 \longrightarrow \mathrm{E} ^ {\prime} \stackrel {f} {\longrightarrow} \mathrm{E} \stackrel {g} {\longrightarrow} \mathrm{E} ^ {\prime \prime}
\]

该序列

\[
0 \longrightarrow \mathrm{T} (\mathrm{E} ^ {\prime}) \xrightarrow {f _ {\mathrm{T}}} \mathrm{T} (\mathrm{E}) \xrightarrow {g _ {\mathrm{T}}} \mathrm{T} (\mathrm{E} ^ {\prime \prime})
\]

是正合的。

\[
\operatorname{Ker} \left(g _ {\mathrm{T}}\right) = \operatorname{Ker} (g) \cap \mathrm{T} (\mathrm{E}) = f \left(\mathrm{E} ^ {\prime}\right) \cap \mathrm{T} (\mathrm{E}) = f \left(\mathrm{T} \left(\mathrm{E} ^ {\prime}\right)\right) = \operatorname{Im} \left(f _ {\mathrm{T}}\right).
\]

命题 25. 设 A 是一个整环，且 \((\mathbf{E}_{\iota})\) 是一族 A-模；则

\begin{enumerate}
\def\labelenumi{(\arabic{enumi})}
\setcounter{enumi}{36}
\tightlist
\item
\end{enumerate}

\[
\mathrm{T} \left(\bigoplus_ {\iota} \mathrm{E} _ {\iota}\right) = \bigoplus_ {\iota} \mathrm{T} (\mathrm{E} _ {\iota}).
\]

设 \((x_{\iota})\) 是 \(\bigoplus_{\iota} E_{\iota}\) 的一个元素，使得 \(x_{\iota} \in T(E_{\iota})\) 对所有 \(\iota\) 成立；则每一个 \(x_{\iota}\) 都被 A 的一个元素 \(\alpha_{\iota} \neq 0\) 所零化，并且可以假设 \(\alpha_{\iota} = 1\) （当 \(x_{\iota} = 0\) 时）；由于族 \((x_{\iota})\) 具有有限支撑，A 的元素 \(\alpha = \prod_{\iota} \alpha_{\iota}\) 是有定义的，并且 \(\neq 0\) ；它显然零化 \(\bigoplus_{\iota} x_{\iota}\) ，因此 \(\bigoplus_{\iota} T(E_{\iota}) \subset T\left(\bigoplus_{\iota} E_{\iota}\right)\) ；其逆命题是显然的。

若 E 和 F 是两个 A-模，则显然 \(T(E \otimes_{A} F)\) 包含 \(T(E) \otimes_{A} F\) 和 \(E \otimes_{A} T(F)\) 的典范像；但可以给出无挠 A-模 E, F 的例子，使得 \(T(E \otimes_{A} F) \neq 0\) （习题 31）。

注意，挠模的无限积不一定是挠模；例如，在 Z-模 \(\prod_{n=1}^{\infty}\left(\mathbf{Z}/p^{n}\mathbf{Z}\right)\) （p 为整数， \textgreater1）中，所有坐标均为 1 的元素是自由的。

命题 26. 设 A 是一个整环，K 是其分式域，E 是一个 A-模，且 \(\mathbf{E}_{(\mathbf{K})} = \mathbf{K} \otimes_{\mathbf{A}} \mathbf{E}\) 是通过扩张算子环得到的 K 上的向量空间；设 \(\phi\) 表示典范 A-线性映射 \(x \mapsto 1 \otimes x\) ，它把 E 映入 \(\mathbf{E}_{(\mathbf{K})}\) .

\begin{enumerate}
\def\labelenumi{(\roman{enumi})}
\item
  \(\mathbf{E}_{(\mathbf{K})}\) 的每一个元素都具有形式 \(\lambda^{-1}\phi(x)\) ，其中 \(\lambda \in \mathbf{A}\) , \(\lambda \neq 0\) ， \(x \in \mathbf{E}\) .
\item
  \(\phi\) 的核是 E 的挠子模 T(E)。
\item
  \(\mathrm{E}_{(\mathbb{K})}\) 的每一个元素都具有形式 \(z = \sum_{i=1}^{n} \xi_i \phi(x_i)\) ，其中 \(\xi_i \in \mathbb{K}\) ， \(x_{i}\in \mathbf{E};\) 对所有 \(i\) ，存在 \(\alpha_{i}\in \mathbf{A}\) ，使得 \(\alpha_{i}\neq 0\) 且 \(\alpha_{i}\xi_{i}\in \mathbf{A}\) ；若 \(\alpha = \prod_{i = 1}^{n}\alpha_{i},\) 则 \(\alpha \neq 0\) 且 \(\alpha \xi_{i} = \beta_{i}\in \mathbf{A}\) 对所有 \(i\) 成立，因此，在 \(\mathrm{E}_{(\mathbf{K})}\) 中，
\end{enumerate}

\[
z = \alpha^ {- 1} (\alpha z) = \alpha^ {- 1} \sum_ {i = 1} ^ {n} \beta_ {i} \phi (x _ {i}) = \alpha^ {- 1} \phi \left(\sum_ {i = 1} ^ {n} \beta_ {i} x _ {i}\right)
\]

因为 \(\phi\) 是 A-线性的。

\begin{enumerate}
\def\labelenumi{(\roman{enumi})}
\setcounter{enumi}{1}
\tightlist
\item
  若 \(x \neq 0\) 在 \(\mathbf{E}\) 中不是自由的，则存在 \(\alpha \neq 0\) 于 \(\mathbf{A}\) 中，使得 \(\alpha x = 0\) ，从而 \(\alpha \phi(x) = \phi(\alpha x) = 0\) 于 \(\mathrm{E}_{(\mathbb{K})}\) 中，这蕴含 \(\phi(x) = 0\) 。反之，假设对某个 \(x \in \mathbf{E}\) 有 \(1 \otimes x = 0\) 于 \(\mathrm{E}_{(\mathbb{K})}\) 中；我们证明 \(x\) 是 \(\mathbf{E}\) 的挠元素。考虑集合 \(\mathfrak{M}\) ，它由 \(\mathbf{K}\) 的单生成子 A-子模组成；在包含关系下它是一个右有向集，因为任意两个元素 \(\alpha, \beta\) （属于 \(\mathbf{K}\) ）可以写成 \(\alpha = \zeta^{-1}\xi, \beta = \zeta^{-1}\eta\) ，其中 \(\xi, \eta, \zeta\) 属于 \(\mathbf{A}\) 且 \(\zeta \neq 0\) ，因此 \(\mathbf{A}. \alpha \subset \mathbf{A}. \zeta^{-1}\) 且 \(\mathbf{A}. \beta \subset \mathbf{A}. \zeta^{-1}\) 。此外， \(\mathbf{K}\) 是诸模 \(\mathbf{M} \in \mathfrak{M}\) 的并，因此可以被视为由诸模 \(\mathbf{M} \in \mathfrak{M}\) 和典范嵌入所定义的正向系统的正向极限（§ 6，第 2 号，注记）。因此在典范同构的意义下， \(\mathrm{E}_{(\mathbb{K})} = \varinjlim (\mathbf{M} \otimes_{\mathbf{A}} \mathbf{E})\) （§ 6，第 3 号，命题 7），且关系 \(1 \otimes x = 0\) 于 \(\mathrm{E}_{(\mathbb{K})}\) 中蕴含存在一个
\end{enumerate}

\(M \in \mathfrak{M}\) ，使得 \(1 \in M\) 且 \(1 \otimes x = 0\) 于张量积 \(M \otimes_{A} E\) 中（集合论，III，§ 7，第 5 号，引理 1）。还可进一步假设（如有必要，把 \(M\) 换成 \(M' \supset M\) ，即 \(K\) 的一个单生成子模） \(M = A. \gamma^{-1}\) ，其中 \(\gamma \in A\) 且 \(\gamma \neq 0\) 。现在映射 \(\xi \mapsto \gamma\xi\) 是 \(M\) 到 A-模 \(A\) 上的一个同构；另一方面，典范同构 \(A \otimes_{A} E \to E\) （§ 3，第 4 号，命题 4）把 \(\xi \otimes x\) 映到元素 \(\xi x\) （它属于 \(E\) ）；因此存在一个同构 \(M \otimes_{A} E \to E\) ，它把张量积 \(\xi \otimes x\) 映到元素 \((\gamma\xi)x\) （属于 \(E\) ）。假设 \(1 \otimes x = 0\) 于 \(M \otimes_{A} E\) 中因此蕴含 \(\gamma x = 0\) .

注记。设 \(\alpha^{-1}\phi(x)\) , \(\beta^{-1}\phi(y)\) 是 \(\mathbf{E}_{(\mathbb{K})}\) 的两个元素，其中 \(\alpha\in\mathbf{A}\) , \(\beta\in\mathbf{A}\) , \(x\in\mathbf{E}\) , \(y\in\mathbf{E}\) , \(\alpha\beta\neq0\) 。要使 \(\alpha^{-1}\phi(x)=\beta^{-1}\phi(y)\) ，必要且充分的条件是 \(\beta x-\alpha y\) 是 \(\mathbf{E}\) 的一个挠元素，因为这个关系等价于 \(\beta\phi(x)=\alpha\phi(y)\) ，后者也可以写成 \(\phi(\beta x-\alpha y)=0\) .

推论 1. 若 E 是无挠 A-模，则典范映射 \(\phi: \mathbf{E} \to \mathbf{E}_{(\mathbb{K})}\) 是单射。

回顾（§ 5，第 1 号）：对于 E 到 K 上向量空间 F 的每个 A-线性映射，存在且仅存在一个 K-线性映射 \(\tilde{f}\colon\mathbf{E}_{(\mathbb{K})}\to\mathbf{F}\) ，使得 \(f=\tilde{f}\circ\phi\) ；我们将称 \(\tilde{f}\) 与 \(f\) 相关联.

推论 2. 设 \(f\) 是把 \(\mathbf{E}\) 映入一个向量空间 \(\mathbf{F}\) （ \(\mathbf{K}\) 上的）的 A-线性映射；若 \(\operatorname{Ker}(f) \subset \mathrm{T}(\mathbf{E})\) ，则 K-线性映射 \(\bar{f}\) （与 \(f\) 相关联的）是单射。

我们把 \(\operatorname{Ker}(\bar{f})\) 的一个元素写成 \(\lambda^{-1}\phi(x)\) 的形式，其中 \(\lambda \in \mathbf{A}\) , \(\lambda \neq 0\) , \(x \in \mathbf{E}\) ；关系 \(\bar{f}(\lambda^{-1}\phi(x)) = 0\) 等价于 \(\lambda^{-1}\bar{f}(\phi(x)) = 0\) 于 \(\mathbf{F}\) 中，从而等价于 \(f(x) = \bar{f}(\phi(x)) = 0\) 。根据假设，这意味着 \(x \in \mathrm{T}(\mathbf{E})\) ，因此 \(\phi(x) = 0\) ，这证明了该推论。

推论 3. 设 E 是一个 A-模，g 是 E 到 K 上向量空间 F 的一个 A-线性映射，使得 \(g(\mathbf{E})\) 生成 F 且 \(\operatorname{Ker}(g) \subset \mathrm{T}(\mathbf{E})\) 。则与 g 相关联的 K-线性映射 \(\bar{g}\) 是 \(\mathrm{E}_{(\mathbb{K})}\) 到 F 上的一个同构。

\(\bar{g}\) 由推论 2 可知是单射的，而 \(g(\mathbf{E})\) 生成 F 这一假设蕴含 \(\bar{g}\) 是满射。

对于每个 A-模 E，向量空间 \(E_{(K)}\) 称为与 E 相伴。对 E 的每个子集 S，S 在 K 上的秩（或按语言习惯的滥用，简称为 S 的秩）是 S 的典范像 \(\phi(S)\) 在 \(E_{(K)}\) 中的秩，换言之（第 2 号，定义 1），即 \(E_{(K)}\) 中由 \(\phi(S)\) 生成的向量子空间在 K 上的维数.

当 E 是无挠 A-模时，通常把它与其典范像 \(\phi(\mathbf{E})\) （在 \(\mathrm{E}_{(\mathbb{K})}\) 中）等同。按照这一约定，E 的每个生成系都包含 \(\mathrm{E}_{(\mathbb{K})}\) 的一组基（第 1 号，定理 2）。特别地：

推论 4. 每个有限生成的 A-模都是有限秩的。

注意，这个推论的逆命题不一定成立；例如，Q 是一个秩为 1 的 Z-模，但它在 Z 上不是有限生成的。

回忆（§ 5，第 1 号）：对于每个线性映射 \(f: \mathbf{E} \to \mathbf{E}'\) （其中 \(\mathbf{E}\) 和 \(\mathbf{E}'\) 是 A-模）， \(f_{(\mathbf{K})}\) 表示 K-线性映射 \(1_{\mathbf{K}} \otimes f: \mathbf{E}_{(\mathbf{K})} \to \mathbf{E}_{(\mathbf{K})}'\) .

命题 27. 对于每一个正合序列

\[
\mathrm{E} ^ {\prime} \xrightarrow {f} \mathrm{E} \xrightarrow {g} \mathrm{E} ^ {\prime \prime}
\]

的A-线性映射，相应的K-线性映射序列

\[
\mathbf {E} _ {(\mathbf {K})} ^ {\prime} \xrightarrow {f _ {(\mathbf {K})}} \mathbf {E} _ {(\mathbf {K})} \xrightarrow {g _ {(\mathbf {K})}} \mathbf {E} _ {(\mathbf {K})} ^ {\prime \prime}
\]

是正合的。

假设 \(g_{(\mathbb{K})}(\lambda^{-1} \otimes x) = 0\) ，其中 \(\lambda \in \mathbf{A}\) , \(\lambda \neq 0\) , \(x \in \mathbb{E}\) ；这等价于 \(\lambda^{-1} \otimes g(x) = 0\) 于 \(\mathrm{E}_{(\mathbb{K})}'\) 中，从而也等价于

\[
1 \otimes g (x) = \lambda (\lambda^ {- 1} \otimes g (x)) = 0;
\]

由命题 26，存在 A 中的 \(\alpha \neq 0\) ，使得 \(\alpha g(x) = 0\) 于 \(\mathbf{E}^{\prime \prime}\) 中，或者也 \(g(\alpha x) = 0\) 。根据假设，因此存在 \(x' \in \mathbf{E}'\) ，使得 \(\alpha x = f(x')\) ，从而 \(\lambda^{-1} \otimes x = f_{(\mathbf{K})}(\alpha^{-1}\lambda^{-1} \otimes x')\) ，这证明了该命题。

推论 1. 若 \(\mathbf{E}'\) 是 \(\mathbf{E}\) 的一个子模，则 \(\mathbf{E}_{(\mathbf{K})}'\) 被典范地等同于 \(\mathbf{E}_{(\mathbf{K})}\) 的一个向量子空间，且 \((\mathbf{E} / \mathbf{E}')_{(\mathbf{K})}\) 被等同于 \(\mathbf{E}_{(\mathbf{K})} / \mathbf{E}_{(\mathbf{K})}'\) .

只需将命题 27 应用于正合序列

\[
0 \rightarrow \mathrm{E} ^ {\prime} \rightarrow \mathrm{E} \rightarrow \mathrm{E} / \mathrm{E} ^ {\prime} \rightarrow 0.
\]

推论 2. 对于每一个 A-线性映射 \(f: \mathrm{E} \to \mathrm{F}\) , \(\operatorname{Ker}(f_{(\mathbf{K})}) = (\operatorname{Ker}(f))_{(\mathbf{K})}\) , \(\operatorname{Im}(f_{(\mathbf{K})}) = (\operatorname{Im}(f))_{(\mathbf{K})}\) , \(\operatorname{Coker}(f_{(\mathbf{K})}) = (\operatorname{Coker}(f))_{(\mathbf{K})}\) 在典范同构的意义下。特别地， \(f_{(\mathbf{K})}\) 为单射（相应地，满射，相应地，零映射）的必要且充分条件是 \(\operatorname{Ker}(f) \subset \mathrm{T}(\mathrm{E})\) （相应地， \(\operatorname{Coker}(f)\) 是一个挠模，相应地， \(\operatorname{Im}(f) \subset \mathrm{T}(\mathrm{F})\) ).

这由推论 1 和命题 26 得出。

推论 3. 设 E 是一个 A-模，且 \((x_{\lambda})_{\lambda \in L}\) 是 E 的一个元素族。要使 \((x_{\lambda})\) 成为一个自由族，必要且充分的条件是：在向量 K-空间 \(\mathbf{E}_{(\mathbb{K})}\) 中，族 \((1 \otimes x_{\lambda})\) 是自由的。

族 \((x_{\lambda})\) 定义了一个 A-线性映射 \(f: \mathbf{A}^{(L)} \to \mathbf{E}\) ，使得 \(f(e_{\lambda}) = x_{\lambda}\) 对所有 \(\lambda \in \mathbf{L}((e_{\lambda})\) （ \(\mathbf{A}^{(L)})\) 的典范基）成立 ，而说 \((x_{\lambda})\) 是自由的就意味着 \(f\) 是单射。只需将推论 2 应用于 \(f\) ，并注意到 \(\mathbf{A}^{(L)}\) 是无挠的（命题 25）。

\section{向量空间中标量域的限制}

在本段中，K 表示一个域，而 \(K'\) 表示 K 的一个子域。在一个集合 V 上，K 上的右（相应地，左）向量空间结构通过标量的限制，定义了 \(K'\) 上的一个右（相应地，左）向量空间结构。

\subsection{K' 结构的定义}

命题 1. 设 V 是 K 上的一个右向量空间，V' 是 V 的一个子集且是 K' 上的向量子空间。以下条件等价：

\begin{enumerate}
\def\labelenumi{(\alph{enumi})}
\item
  K-线性映射 \(\lambda\) （由 \(\mathrm{V}_{(\mathbb{K})}^{\prime} = \mathrm{V}^{\prime}\otimes_{\mathbb{K}^{\prime}}\mathbf{K}\) 到 V 中的），使得 \(\lambda (x^{\prime}\otimes \xi) = x^{\prime}\xi\) 对所有 \(x^{\prime}\in \mathrm{V}^{\prime},\xi \in \mathbf{K}\) 成立，是双射的。
\item
  每个从 V' 到向量 K-空间 W 中的 K'-线性映射 \(f'\) 都可以唯一地延拓为从 V 到 W 中的一个 K-线性映射 f。
\item
  \(\mathbf{V}'\) 在 \(\mathbf{K}'\) 上的每一组基都是 \(\mathbf{V}\) 在 \(\mathbf{K}\) 上的一组基。
\item
  存在 \(\mathbf{V}'\) 在 \(\mathbf{K}'\) 上的一组基，它同时也是 \(\mathbf{V}\) 在 \(\mathbf{K}\) 上的一组基。
\item
  向量 K-空间 V 由 V' 生成，且 V' 的每个在 K' 上自由的子集在 K 上也是自由的。
\end{enumerate}

我们知道（§ 5，第 1 号）， \(\mathbf{V}_{(\mathbf{K})}^{\prime}\) 具有一个右向量 K-空间结构，在此结构下 \((x^{\prime} \otimes \xi)\eta = x^{\prime} \otimes (\xi\eta)\) （ \(\xi, \eta\) 属于 K ， \(x^{\prime} \in \mathbf{V}^{\prime}\) ），并且对于每个从 V' 到向量 K-空间 W 中的 K'-线性映射 \(f^{\prime}\) ，存在且仅存在一个 K-线性映射 \(\bar{f}^{\prime}\) （从 \(\mathbf{V}_{(\mathbf{K})}^{\prime}\) 到 W 中的），使得 \(\bar{f}^{\prime}(x^{\prime} \otimes 1) = f^{\prime}(x^{\prime})\) 对所有 \(x^{\prime} \in \mathbf{V}^{\prime}\) 成立。若 \(j\) 是 V' 到 V 中的典范单射，则 \(\lambda\) 正是相应的 K-线性映射 \(\bar{j}\) 。若 \(\lambda\) 是双射，则对每个 K'-线性映射 \(f^{\prime}: \mathbf{V}^{\prime} \to \mathbf{W}, \bar{f}^{\prime} \circ \lambda^{-1}\) 是 V 到 W 中的延拓 \(f^{\prime}\) 的唯一 K-线性映射；换言之，(a) 蕴含 (b)。反之，若 (b) 成立，则特别地存在一个 K-线性映射 \(\mu\) （从 V 到 \(\mathbf{V}_{(\mathbf{K})}^{\prime}\) 中的），使得 \(\mu(x^{\prime}) = x^{\prime} \otimes 1\) 对所有 \(x^{\prime} \in \mathbf{V}^{\prime}\) 成立；立即可得 \(\mu \circ \lambda = 1_{\mathbf{V}_{(\mathbf{K})}^{\prime}}\) ；另一方面， \(\lambda(\mu(x^{\prime})) = x^{\prime}\) 对所有 \(x^{\prime} \in \mathbf{V}^{\prime}\) 成立，且根据假设 \(j: \mathbf{V}^{\prime} \to \mathbf{V}\) 可以唯一地延拓为 V 的一个自同态，故必然有 \(\lambda \circ \mu = 1_{\mathbf{V}}\) ，这就完成了 (a) 与 (b) 等价的证明。

对每一组基 \(\mathbf{B}'\) （ \(\mathbf{V}'\) 在 \(\mathbf{K}'\) 上的），集合 \(\mathbf{B}\) 的元素是形如 \(v' \otimes 1\) 的元素（属于 \(\mathbf{V}_{(\mathbf{K})}'\) ，其中 \(v'\) 遍历 \(\mathbf{B}'\) ），它是 \(\mathbf{V}_{(\mathbf{K})}'\) 在 \(\mathbf{K}\) 上的一组基（ \(\S 5\) ，第 1 号，命题 4），且 \(\lambda(\mathbf{B}) = \mathbf{B}'\) 。要使 \(\lambda\) 是双射， \(\lambda\) 必须把 \(\mathbf{V}_{(\mathbf{K})}'\) 的每组基映为 \(\mathbf{V}\) 在 \(\mathbf{K}\) 上的一组基，而且只需对 \(\mathbf{V}_{(\mathbf{K})}'\) 的单独一组基如此即可（ \(\S 1\) ，第 11 号，命题 17 的推论 2）。这证明了 (a)、(c) 和 (d) 的等价性。

由于 \(\mathbf{V}'\) 的每个在 \(\mathbf{K}'\) 上自由的子集都包含在 \(\mathbf{V}'\) 在 \(\mathbf{K}'\) 上的一组基之中（§ 7，第 1 号，定理 2），(c) 蕴含 (e)。最后，假设 (e) 成立；若 \(\mathbf{B}'\) 是 \(\mathbf{V}'\) 在 \(\mathbf{K}'\) 上的一组基，则它是 \(\mathbf{V}\) 在 \(\mathbf{K}\) 上的一个自由子集；另一方面， \(\mathbf{B}'\) 生成 \(\mathbf{V}'\) 于 \(\mathbf{K}'\) 上，故根据假设它也生成 \(\mathbf{V}\) 于 \(\mathbf{K}\) 上；因此 \(\mathbf{B}'\) 是 \(\mathbf{V}\) 在 \(\mathbf{K}\) 上的一组基，这证明了 (e) 蕴含 (c)。

定义 1. 设 V 是域 K 上的一个右向量空间，K' 是 K 的一个子域。V 的每个满足命题 1 的等价条件的向量子 K'-空间 V' 称为 V 上的一个 K'-结构。

例. 设 B 是 V 在 K 上的一组基。对 K 的每个子域 \(K'\) ，由 B 生成的 V 的向量子 \(K'\) -空间以 B 作为它在 \(K'\) 上的一组基，因而是 V 上的一个 \(K'\) -结构。*例如，若 K 是交换的且 V 取为多项式 K-代数 \(K[X_{1},\ldots,X_{n}]\) ，那么对 K 的每个子域 \(K'\) ， \(K'[X_{1},\ldots,X_{n}]\) 是一个 \(K'\) -结构（在 \(V\) 上）。*

\subsection{子空间的有理性}

定义 2. 设 V 是 K 上的一个右向量空间，带有一个 K'-结构 V'。V 中的一个向量称为在 K' 上有理的，如果它属于 V'。V 的一个向量子 K-空间 W 称为在 K' 上有理的，如果它（在 K 上）由在 K' 上有理的向量生成。

设 \((v_{\iota}^{\prime})_{\iota\in I}\) 是 \(V^{\prime}\) 在 \(K^{\prime}\) 上的一组基，因而也是 \(V\) 在 \(K\) 上的一组基（第 1 号，命题 1）。要使向量 \(x = \sum_{\iota} v_{\iota}^{\prime} \xi_{\iota}\) （属于 \(V\) ）在 \(K^{\prime}\) 上有理，必要且充分的条件是 \(\xi_{\iota} \in K^{\prime}\) 对所有 \(\iota \in I\) 成立。

若 W 是 V 的一个在 \(K'\) 上有理的向量子 K-空间，则由定义 2 可知， \(W' = W \cap V'\) 是 W 的一个向量子 \(K'\) -空间，它在 K 上生成 W；另一方面， \(W'\) 的每个在 \(K'\) 上自由的子集在 K 上也是自由的，因为它包含于 \(V'\) （第 1 号，命题 1）。由此得出（第 1 号，命题 1）： \(W'\) 是 W 上的一个 \(K'\) -结构，称为由 V 上的 \(K'\) -结构 \(V'\) 诱导的结构。

对 V' 的每个向量子 K'-空间 W'，我们将用 W'.K 表示 V 中由 W' 的元素以 K 中系数作线性组合所构成的向量子 K-空间。

命题 2. 设 V 是 K 上的一个右向量空间，V' 是 V 上的一个 K'-结构。映射 W' → W'.K 是从 V' 的向量子 K'-空间集合到 V 的在 K' 上有理的向量子 K-空间集合上的一个双射，其逆双射为 W → W ∩ V'。

显然，双射 \(\lambda^{-1}: V \to V' \otimes_{K'} K\) （即第 1 号命题 1 中定义的双射 \(\lambda\) 的逆映射）把每个向量子 \(K'\) -空间 \(W'\) （ \(V'\) 的）映到它在典范单射 \(x' \mapsto x' \otimes 1\) 下的像上，并把 \(W'. K\) 映到 \(W' \otimes_{K'} K\) 上；因此命题 2 的断言是定义 2 和 § 7 第 9 号命题 19 的推论。

推论 1. V 的在 K' 上有理的向量子 K-空间的每一组和与每一个交都是在 K' 上有理的子空间。

关于和的断言是显然的。另一方面，若 \((\mathbf{W}_{\iota}')_{\iota \in I}\) 是向量子 \(\mathbf{K}'\) -空间的一个族（在 \(\mathbf{V}'\) 中），则 \(\left(\bigcap_{\iota \in I} \mathbf{W}_{\iota}'\right) \otimes_{\mathbb{K}'} \mathbf{K} = \bigcap_{\iota \in I} (\mathbf{W}_{\iota}' \otimes_{\mathbb{K}'} \mathbf{K})\) （§ 7，第 7 号，命题 14 的推论），这就证明了该推论。

V 在 K 上的一组基 B 称为在 \(K'\) 上有理的，如果它由在 \(K'\) 上有理的向量组成。

推论 2. V 在 K 上的每组在 \(\mathbf{K}'\) 上有理的基都是 \(\mathbf{V}'\) 在 \(\mathbf{K}'\) 上的一组基。

若 \(W'\) 是由 B 生成的向量子 \(K'\) -空间（ \(V'\) 中的），则 \(W'. K = V = V'. K\) ，于是根据命题 2 有 \(V' = W'\) 。

\subsection{线性映射的有理性}

定义 3. 设 \(\mathbf{V}_1, \mathbf{V}_2\) 是 \(\mathbf{K}\) 上的两个右向量空间，分别带有 \(\mathbf{K}'\) -结构 \(\mathbf{V}_1'\) , \(\mathbf{V}_2'\) 。一个 K-线性映射 \(f: \mathbf{V}_1 \to \mathbf{V}_2\) 称为在 \(\mathbf{K}'\) 上有理的，如果 \(f(\mathbf{V}_1') \subset \mathbf{V}_2'\) 。

若 \(V_{3}\) 是 K 上的第三个右向量空间，带有一个 \(K^{\prime}\) -结构 \(V_{3}^{\prime}\) ，且 K-线性映射 \(g:V_{2}\rightarrow V_{3}\) 在 \(K^{\prime}\) 上有理，则显然 \(g\circ f:V_{1}\rightarrow V_{3}\) 在 \(K^{\prime}\) 上有理。

命题 3. 设 \(\mathbf{V}_1, \mathbf{V}_2\) 是 \(\mathbf{K}\) 上的两个右向量空间，且 \(\mathbf{V}_1', \mathbf{V}_2' \mathbf{K}'\) -结构分别给定在 \(\mathbf{V}_1, \mathbf{V}_2\) 上。 \(\mathbf{V}_1\) （相应地 \(\mathbf{V}_2\) ）典范地等同于 \(\mathbf{V}_1' \otimes_{\mathbb{K}'} \mathbf{K}\) （相应地 \(\mathbf{V}_2' \otimes_{\mathbb{K}'} \mathbf{K}\) ）（第 1 号，命题 1）。

\begin{enumerate}
\def\labelenumi{(\roman{enumi})}
\item
  映射 \(f' \mapsto f' \otimes 1_{\mathbf{K}} = f_{(\mathbf{K})}'\) 是从 \(\operatorname{Hom}_{\mathbf{K}'}(\mathbf{V}_1', \mathbf{V}_2')\) 到 \(\mathbf{V}_1\) 到 \(\mathbf{V}_2\) 中的在 \(\mathbf{K}'\) 上有理的 K-线性映射集合上的一个双射；逆双射把每个 K-线性映射 \(f: \mathbf{V}_1 \to \mathbf{V}_2\) （在 \(\mathbf{K}'\) 上有理的）对应到 K'-线性映射 \(f': \mathbf{V}_1' \to \mathbf{V}_2'\) ，它与 f 到 \(\mathbf{V}_1'\) 上的限制具有相同的图像。
\item
  对每个 K-线性映射 \(f: \mathbf{V}_1 \to \mathbf{V}_2\) （在 \(\mathbf{K}'\) 上有理的），
\end{enumerate}

\[
f (\mathrm{V} _ {1} ^ {\prime}) = f (\mathrm{V} _ {1}) \cap \mathrm{V} _ {2} ^ {\prime} \quad and \quad \bar {f} ^ {- 1} (\mathrm{V} _ {2} ^ {\prime}) = \mathrm{V} _ {1} ^ {\prime} + \operatorname{Ker} (f).
\]

\begin{enumerate}
\def\labelenumi{(\roman{enumi})}
\item
  显然，在上述等同下，若 \(f': \mathbf{V}_1' \to \mathbf{V}_2'\) 是一个 K'-线性映射，则 \(f_{(\mathbf{K})}' = f' \otimes 1_{\mathbf{K}}\) 在 \(\mathbf{K}'\) 上有理，且 \(f'\) 是与 \(f_{(\mathbf{K})}'\) 到 \(\mathbf{V}_1'\) 上的限制具有相同图像的映射。反之，若 \(f: \mathbf{V}_1 \to \mathbf{V}_2\) 是一个在 \(\mathbf{K}'\) 上有理的 K-线性映射，且 \(f': \mathbf{V}_1' \to \mathbf{V}_2'\) 与 \(f\) 到 \(\mathbf{V}_1'\) 上的限制具有相同的图像，则 \(f\) 和 \(f_{(\mathbf{K})}'\) 在 \(\mathbf{V}_1'\) 上一致，后者是 \(\mathbf{V}_1\) 在 \(\mathbf{K}\) 上的一个生成系，故 \(f = f_{(\mathbf{K})}'\) 。
\item
  若 \(f = f' \otimes 1_{\mathbb{K}}\) ，则 \(f(\mathbf{V}_1) = f(\mathbf{V}_1' \otimes_{\mathbb{K}'} \mathbf{K}) = f'(\mathbf{V}_1') \otimes_{\mathbb{K}'} \mathbf{K}\) ，且由于 \(f'(\mathbf{V}_1') \subset \mathbf{V}_2'\) ，有 \(f(\mathbf{V}_1') = f'(\mathbf{V}_1') = f(\mathbf{V}_1) \cap \mathbf{V}_2'\) （§ 7，第 9 号，命题 19）；公式 \(f^{-1}(\mathbf{V}_2') = \mathbf{V}_1' + \operatorname{Ker}(f)\) 立即得出。
\end{enumerate}

推论 1. 在命题 3 的记号下， \(\operatorname{Im}(f) = (\operatorname{Im}(f'))_{(\mathbb{K})}\) ，

\[
\operatorname{Ker} (f) = (\operatorname{Ker} (f ^ {\prime})) _ {(\mathbb {K})}, \quad \operatorname{Coker} (f) = (\operatorname{Coker} (f ^ {\prime})) _ {(\mathbb {K})}.
\]

特别地，要使 \(f\) 是单射（相应地，满射、零映射），必要且充分的条件是 \(f'\) 如此。若 \(f\) 是双射的，则其逆映射在 \(\mathbf{K}'\) 上有理。

这是 § 7 第 9 号命题 18 的推论的一个特例。

推论 2. 设 \(f: \mathbf{V}_1 \to \mathbf{V}_2\) 是一个在 \(\mathbf{K}'\) 上有理的 K-线性映射。对每个向量子 K-空间 \(\mathbf{W}_1\) （ \(\mathbf{V}_1\) 的；相应地 \(\mathbf{W}_2\) ， \(\mathbf{V}_2\) 的），若它在 \(\mathbf{K}'\) 上有理，则 \(f(\mathbf{W}_1)\) （相应地 \(f(\mathbf{W}_2)\) ）是 \(\mathbf{V}_2\) （相应地 \(\mathbf{V}_1\) ）的在 \(\mathbf{K}'\) 上有理的向量子 K-空间。

在命题 3 的记号下，对每个向量子 K'-空间 W \(_{1}\) （属于 V \(_{1}\) ），有 f \(_{(K)}\) (W \(_{1}\) ' ⊗ \(_{K'}\) K) = f'(W \(_{1}\) ) ⊗ \(_{K'}\) K；由此得出关于 W \(_{1}\) 的断言（第 2 号，命题 2）。另一方面，设 W \(_{2}\) 是一个向量子 K'-空间（ V \(_{2}\) 的），并令 g' 为典范 K'-线性映射 V \(_{2}\) → V \(_{2}\) /W \(_{2}\) ；于是 \(f'^{-1}(W_{2}') = \text{Ker}(g' \circ f')\) ；因此，由推论 1，

\[
\bar {f} _ {(\mathbf {K})} ^ {\prime - 1} (\mathbf {W} _ {2} ^ {\prime} \otimes_ {\mathbf {K} ^ {\prime}} \mathbf {K}) = \bar {f} ^ {\prime - 1} (\mathbf {W} _ {2} ^ {\prime}) \otimes_ {\mathbf {K} ^ {\prime}} \mathbf {K},
\]

由此得出关于 \(W_{2}\) 的断言。

设 \(V_{1}, V_{2}\) 是两个右向量 K-空间，分别带有 \(K'\) -结构 \(V_{1}', V_{2}'\) 。显然， \(V_{1}' \times V_{2}'\) 是一个 \(K'\) -结构（在 \(V_{1} \times V_{2}\) 上的），称为 \(K'\) -结构 \(V_{1}'\) 与 \(V_{2}'\) 的积。

命题 4. 要使 K-线性映射 \(f: \mathbf{V}_1 \to \mathbf{V}_2\) 在 \(\mathbf{K}'\) 上有理，必要且充分的条件是它的图像 \(\Gamma\) 在 \(\mathbf{K}'\) 上有理（对于 \(\mathbf{K}'\) -积结构，即 \(\mathbf{V}_1 \times \mathbf{V}_2\) 上的那个而言）。

设 \(g\) 为映射 \(x_1 \mapsto (x_1, f(x_1))\) （从 \(V_1\) 到 \(V_1 \times V_2\) 中的）；这是一个使得 \(\Gamma = g(V_1)\) 的 K-线性映射；若 \(f\) 在 \(K'\) 上有理，则根据命题 3 的推论 2， \(\Gamma\) 亦然。反之，设 \(\Gamma\) 在 \(K'\) 上有理，并令它具有一个 \(K'\) -结构（由 \(V_1 \times V_2\) 上的结构诱导的）；由定义立即得出，限制 \(p_1, p_2\) （到 \(\Gamma\) 上的、投影 \(pr_1, pr_2\) 的）是 K-线性映射，在 \(K'\) 上有理，分别从 \(\Gamma\) 到 \(V_1\) 和 \(V_2\) 中。由于 \(p_1\) 是双射的，其逆映射 \(q_1\) 在 \(K'\) 上有理（命题 3 的推论 1），因而 \(f = p_2 \circ q_1\) 亦然。

\subsection{有理线性形式}

设 V 是 K 上的一个右向量空间，带有一个 \(K^{\prime}\) -结构 \(V^{\prime}\) 。由于 \(K_{d}^{\prime}\) 是一个 \(K^{\prime}\) -结构（在右向量 K-空间 \(K_{d}\) 上的），我们可以把线性形式 \(x^{*} \in V^{*}\) （在 \(K^{\prime}\) 上有理的）定义为 V 到 \(K_{d}\) 中的线性映射（在 \(K^{\prime}\) 上有理的，相对于

上的 K'-结构，即 V 上的和 \(K_{d}\) 上的）。根据第 3 号命题 3，这些线性形式的集合 \(R'\) 是对偶空间 \(V'^*\) （ \(V'\) 的对偶）在复合映射

\[
\mathrm{V} ^ {\prime *} \xrightarrow {\phi} \mathrm{K} \otimes_ {\mathrm{K} ^ {\prime}} \mathrm{V} ^ {\prime *} \xrightarrow {\upsilon} \mathrm{V} ^ {*}\tag{1}
\]

下的像，其中 \(\phi(x^{\prime*}) = 1 \otimes x^{\prime*}\) ，而 \(\upsilon(\xi \otimes x^{\prime*})\) 是 V 上的线性形式 \(y^{*}\) ，使得 \(y^{*}(x^{\prime}) = \xi\langle x^{\prime*}, x^{\prime} \rangle\) 对所有 \(x^{\prime} \in V^{\prime}\) 成立（§ 5，第 4 号）。我们知道这个映射是单射（§ 7，第 9 号，命题 20 和 19），且显然 \(R^{\prime}\) 是一个左向量子 \(K^{\prime}\) -空间（ \(V^{*}\) 的）；此外， \(R^{\prime}\) 的每个在 \(K^{\prime}\) 上自由的子集在 K 上也是自由的。但一般说来， \(R^{\prime}\) 不一定在 K 上生成 \(V^{*}\) ，因而并不定义一个 \(K^{\prime}\) -结构于 \(V^{*}\) 之上（习题 2）。然而，若 V 在 K 上是有限维 \(n\) 的，则 \(V^{\prime*}\) 是维数 \(n\) 的（在 \(K^{\prime}\) 上），于是 \(R^{\prime}\) 典范地定义一个 \(K^{\prime}\) -结构于 \(V^{*}\) 之上。

命题 5. 设 V 是 K 上的右向量空间，V' 是 V 上的一个 K'-结构，W 是 V 的一个向量子 K-空间。要使 W 在 K' 上有理，必要且充分的条件是：存在一个由在 K' 上有理的线性形式组成的集合 H ⊂ V*，使得 W 是 H 在 V 中的正交补（§ 2，第 4 号）。

设 H 是 V* 的一个子集，其元素是在 K' 上有理的线性形式。对所有 \(x^{*} \in H\) ， \(x^{*}\) 的核是 V 的一个在 K' 上有理的向量子 K-空间（第 3 号，命题 3 的推论 2）；因此这些核的交也是 V 的一个在 K' 上有理的向量子 K-空间（第 2 号，命题 2 的推论 1）。

反之，设 W 是 V 的一个在 \(K'\) 上有理的向量子 K-空间，于是 W 等同于 \(W' \otimes_{K'} K\) ，其中 \(W' = W \cap V'\) （第 2 号，命题 2）。要使线性形式 \(x'^* \in V'^*\) 在 \(W'\) 上为零，必要且充分的条件是它在 (1) 下对应的线性形式 \(x^* \in V^*\) 在 W 上为零，因为根据第 3 号命题 3 的推论 1，

\[
\operatorname{Ker} (x ^ {*}) = \left(\operatorname{Ker} (x ^ {\prime *})\right) \otimes_ {\mathbb {K} ^ {\prime}} \mathrm{K} \quad \text { and } \quad \operatorname{Ker} (x ^ {* ^ {\prime}}) = (\operatorname{Ker} (x ^ {*})) \cap \mathrm{V} ^ {\prime}.
\]

设 \(H'\) 是 \(W'\) 在 \(V'^*\) 中的正交补；我们知道（§ 7，第 5 号，定理 7）， \(W'\) 是 \(H'\) 在 \(V'\) 中的正交补；若 \(H\) 是 \(H'\) 在映射 (1) 下在 \(V*\) 中的像，则由上述可知， \(W\) 是 \(H\) 在 \(V\) 中的正交补，这里用到第 7 号命题 14 的推论。

\subsection{对线性系统的应用}

命题 6. (i) 给定一个齐次线性方程组

\[
\sum_ {\iota \in I} \alpha_ {\mu \iota} \xi_ {\iota} = 0 (\mu \in M)\tag{2}
\]

其系数 \(\alpha_{\mu \iota}\) 属于 \(K'\) ，该方程组的每个由 K 中元素组成的解（ \(\xi_{\iota}\) ）都是 (2) 的解（ \(\xi_{\iota}'\) ，由 \(K'\) 中元素组成的）的、以 K 中系数作出的线性组合。

\begin{enumerate}
\def\labelenumi{(\roman{enumi})}
\setcounter{enumi}{1}
\tightlist
\item
  给定一个线性方程组
\end{enumerate}

\[
\sum_ {\iota \in I} \alpha_ {\mu \iota} \xi_ {\iota} = \beta_ {\mu} \quad (\mu \in \mathbf {M})\tag{3}
\]

其系数 \(\alpha_{\mu \iota}\) 和右端项 \(\beta_{\mu}\) 属于 \(K'\) ，如果存在该方程组的一个由 K 中元素组成的解，则也存在一个由 \(K'\) 中元素组成的解。

\begin{enumerate}
\def\labelenumi{(\roman{enumi})}
\item
  对每个集合 S，设右向量 K-空间 \(\mathrm{K}_{d}^{(\mathrm{S})}\) 被赋予 \(K^{\prime}\) -结构 \(\mathrm{K}_{d}^{\prime(\mathrm{S})}\) 。设 f 是把 \(\mathrm{K}_{d}^{(1)}\) 映入 \(\mathrm{K}_{d}^{(\mathrm{M})}\) 的 K-线性映射，它把每个向量 \((\xi_{\iota})_{\iota\in\mathbb{I}}\) 映为向量 \((\zeta_{\mu})_{\mu\in\mathbb{M}}\) ，其中 \(\zeta_{\mu}=\sum_{\iota\in\mathbb{I}}\alpha_{\mu\iota}\xi_{\iota}\) 对所有 \(\mu\in M\) 成立。显然 f 在 \(K^{\prime}\) 上有理；它的核 V（即 K 中方程组 (2) 的解集）是 \(\mathrm{K}_{d}^{(\mathrm{I})}\) 的一个在 \(K^{\prime}\) 上有理的子空间（第 3 号，命题 3 的推论 2），因而由 (2) 的在 \(K^{\prime}\) 中的解生成。
\item
  我们把 K 视为一个左向量 \(\mathbf{K}'\) -空间；存在一个 \(\mathbf{K}'\) -线性投影算子 \(p\) ，它把 K 映到其向量子空间 \(\mathbf{K}_s'\) 上（§ 7，第 3 号，命题 4）；若 \((\xi_{\iota})\) 是 (3) 在 K 中的一个解，则 \(\sum_{\iota \in I} \alpha_{\mu \iota} p(\xi_{\iota}) = p\left(\sum_{\iota \in I} \alpha_{\mu \iota} \xi_{\iota}\right) = p(\beta_\mu) = \beta_\mu\) ，这就证明了 \((p(\xi_{\iota}))\) 是 (3) 在 \(\mathbf{K}'\) 中的一个解。
\end{enumerate}

一个环 K 称为在子环 \(K'\) 上（左）忠实平坦的，如果命题 6 对 K 和 \(K'\) 成立；我们将在后面详细研究这一概念（交换代数，I，§3）。

\subsection{最小有理性域}

设 V 是带有一个 K'-结构 V' 的右向量 K-空间。对每个满足 \(K' \subset L \subset K\) 的域 L，我们记 \(V_{L} = V'. L\) ；显然， \(V'\) 在 \(K'\) 上的每组基既是 V 在 K 上的基，也是 \(V_{L}\) 在 L 上的基。因此 \(V_{L}\) 是 V 上的一个 L-结构，而 \(V'\) 是 \(V_{L}\) 上的一个 K'-结构。

命题 7. (i) 设 V 是带有一个 K'-结构 V' 的右向量 K-空间。对每个向量 \(x \in V\) （相应地，V 的每个向量子 K-空间 W），K 的包含 K' 且使得 x（相应地，W）在其上有理的子域 L 的集合有一个最小元 \(\mathrm{K}'(x)\) （相应地 \(\mathrm{K}'(\mathrm{W})\) ）。

\begin{enumerate}
\def\labelenumi{(\roman{enumi})}
\setcounter{enumi}{1}
\tightlist
\item
  设 \(V_{1}\) , \(V_{2}\) 是两个右向量 K-空间，分别带有 \(K^{\prime}\) -结构 \(V_{1}^{\prime}\) , \(V_{2}^{\prime}\) 。对每个从 \(V_{1}\) 到 \(V_{2}\) 中的 K-线性映射 f，K 的包含 \(K^{\prime}\) 且使得 f 在其上有理的子域 L 的集合有一个最小元 \(\mathrm{K}^{\prime}(f)\) 。
\end{enumerate}

我们首先对向量 \(x \in V\) 证明断言 (i)。设 \(B\) 是 \(V\) 的一个在 \(K'\) 上有理的基； \(B\) 是 \(V'\) 在 \(K'\) 上的一组基，并且是 \(V_L\) 在 \(L\) 上的一组基，这里 \(L\) 是任意满足 \(K' \subset L \subset K\) 的域；要使 \(x = \sum_{b \in B} b\xi_b\) 在 \(L\) 上有理，必要且充分的条件是诸 \(\xi_{b}\) 属于 L（第 2 号），因此具有这一性质的最小域 L 是由 K' 与诸 \(\xi_{b}\) （ \(b \in \mathbf{B}\) ）生成的域。

我们接下来证明 (ii)。设 \(\mathbf{B}_1, \mathbf{B}_2\) 分别是 \(\mathbf{V}_1, \mathbf{V}_2\) 的在 \(\mathbf{K}'\) 上有理的基，并对所有 \(b_1 \in \mathbf{B}_1\) 写出 \(f(b_1) = \sum_{b_2 \in \mathbf{B}_2} b_2 \alpha_{b_2 b_1}\) （*族 \((\alpha_{b_2 b_1})\) 正是 \(f\) 关于基 \(\mathbf{B}_1\) 和 \(\mathbf{B}_2\) 的矩阵；参见 § 10，第 4 号*）。由于 \(\mathbf{B}_1\) （相应地 \(\mathbf{B}_2\) ）是 \((\mathbf{V}_1)_L\) （相应地 \((\mathbf{V}_2)_L\) ）在 \(L\) 上的一组基（对每个满足 \(\mathbf{K}' \subset L \subset K\) 的域而言），故要使 \(f\) 在 \(L\) 上有理，必要且充分的条件是诸 \(\alpha_{b_2 b_1}\) 属于 \(L\) ；因此具有这一性质的最小域是由 \(\mathbf{K}'\) 与诸 \(\alpha_{b_2 b_1}\) （ \(b_1 \in \mathbf{B}_1\) , \(b_2 \in \mathbf{B}_2\) ）生成的域。

最后，为了对V的子空间W建立断言(i)，我们首先证明以下引理：

引理 1. 设 V 是带有一个 K'-结构 V' 的右向量 K-空间，W 是 V 的一个向量子 K-空间。则存在 V 的两个在 K' 上有理的向量子 K-空间 W₁、W₂，使得 V 是 W₁ 与 W₂ 的直和，并且若将 V 等同于 W₁ × W₂，则 W 是从 W₁ 到 W₂ 中的一个 K-线性映射 g 的图像。

设 B 为 V 的一个在 \(K'\) 上有理的基。将 § 7 第 1 号定理 2 应用于 W 在 K 上的一组基（把它视为 V 的一个自由子集）以及由该自由子集与 B 的并组成的生成系，可以看出存在 B 的一个子集 C，使得 V 是 W 与由 C 生成的 V 的子空间 \(W_{2}\) 的直和。又设 \(W_{1}\) 为 V 中由 B - C 生成的子空间。由于 \(B \subset V'\) ，显然 \(W_{1}\) 和 \(W_{2}\) 在 \(K'\) 上有理。此外，对所有 \(x \in W_{1}\) ，存在且仅存在一个向量 \(g(x)\) （属于 \(W_{2}\) ），使得 \(x + g(x) \in W\) ，因为 V 是 W 与 \(W_{2}\) 的直和；于是 W 是 g 的图像，且由于 W 是 V 的向量子 K-空间，g 是 K-线性的。

证明了这个引理之后，我们知道 W 在 K 的包含 \(K'\) 的子域 L 上有理当且仅当 g 在 L 上有理（第 3 号，命题 4）。域 \(\mathrm{K}'(g)\) 是使得 g 在 \(\mathrm{K}'(g)\) 上有理的最小域，因而也是使 W 在其上有理的最小域。
